\documentclass[11pt]{article}
\usepackage[letterpaper,margin=1in]{geometry}
\usepackage[T1]{fontenc}
\usepackage{lmodern,microtype}
\usepackage{amsmath,amssymb,amsthm,mathtools,mathrsfs}
\usepackage{enumitem,booktabs,array,xcolor}
\usepackage[colorlinks=true,linkcolor=blue!40!black,citecolor=blue!40!black,urlcolor=blue!40!black]{hyperref}
\usepackage{fancyhdr}
\pagestyle{fancy}
\fancyhf{}
\fancyhead[L]{\small Chiral curvature and metric evolution}
\fancyhead[R]{\small Working draft}
\fancyfoot[C]{\thepage}
\setlength{\headheight}{14pt}
\setlength{\emergencystretch}{2em}
\setcounter{tocdepth}{2}
\numberwithin{equation}{section}
\newtheorem{theorem}{Theorem}[section]
\newtheorem{proposition}[theorem]{Proposition}
\newtheorem{lemma}[theorem]{Lemma}
\newtheorem{corollary}[theorem]{Corollary}
\theoremstyle{definition}
\newtheorem{definition}[theorem]{Definition}
\theoremstyle{remark}
\newtheorem{remark}[theorem]{Remark}
\DeclareMathOperator{\Ric}{Ric}
\DeclareMathOperator{\tr}{tr}
\DeclareMathOperator{\End}{End}
\DeclareMathOperator{\Sym}{Sym}
\DeclareMathOperator{\Vol}{Vol}
\newcommand{\R}{\mathbb R}
\newcommand{\CP}{\mathbb{CP}}
\newcommand{\dd}{\mathrm d}
\newcommand{\id}{\operatorname{id}}

\title{Chiral curvature and metric evolution\\near the round four-sphere}
\author{Edward B. Baker III}
\date{Working draft --- 22 November 2026}
% Background revision: 21 November 2026 (America/New_York), for Edward B. Baker III,
% with OpenAI assistance. Actual model/effort not independently exposed.
% Inputs, passage checks and editorial scope: ../reviews/BACKGROUND_REVISION_20261121.md.
% This revision does not constitute the separate audit of the research claims.
% Introduction progression revised 21 November 2026 (America/New_York),
% for Edward B. Baker III with OpenAI assistance; actual model/effort unavailable.
% Input versions and editorial scope: ../reviews/INTRODUCTION_REWRITE_20261121.md.
% Citation style and paragraph structure revised 22 November 2026 (America/New_York),
% for Edward B. Baker III with OpenAI assistance; actual model/effort unavailable.
% Input version and editorial scope: SP4-CME-DOC-E03 in ../research-progress/REGISTER.md.

\begin{document}
\maketitle
\thispagestyle{plain}
\begin{abstract}
Motivated by the smooth four-dimensional Poincar\'e conjecture, we develop a self-contained comparison between chiral connection evolution and Riemannian metric evolution near the round four-sphere. A definite connection determines a conformal metric through its curvature span and a representative through a curvature volume form. We derive the differential of this reconstruction, its exact velocity for an arbitrary connection path, and the first nonlinear corrections when the connection is induced by a source metric. For Fine's definite-connection gradient flow, the linearized normalized curvature defect satisfies a closed heat equation. A contracted Bianchi identity shows that the flow preserves the image of transverse-traceless metric variations under the linearized chiral Levi--Civita map; the reconstructed metric then obeys four times the linearization of volume-normalized Ricci flow. The same identities give a direct factorization of the round Weyl-energy Hessian and distinguish the second-order connection evolution from the fourth-order metric gradient of the pulled-back energy. We compute the cubic discrepancy between the metric energies, identify the Levi--Civita compatibility defect of the reconstructed geometry, and exhibit finite energy density at the boundary of metric nondegeneracy. The final discussion separates these local results from the existence, preservation, and compactness estimates required for a global application to smooth homotopy four-spheres.
\end{abstract}
\tableofcontents
\newpage



\section{Introduction}\label{sec:introduction}

Isothermal coordinates provide a simple model for using geometry to
recognize the smooth structure of a sphere. Every smooth Riemannian metric
on an oriented real surface can be written locally as
\[
 g=e^{2u}(dx^2+dy^2)=e^{2u}|dz|^2,\qquad z=x+iy.
\]
Orientation-preserving transitions between these charts are holomorphic,
so the metric determines a Riemann surface. If the underlying smooth
surface is closed and homotopy equivalent to \(S^2\), it is orientable
and simply connected. Uniformization then identifies this compact Riemann
surface biholomorphically with \(\CP^1\); the resulting map is a
diffeomorphism to the standard smooth two-sphere~\cite{Donaldson2011}.
This gives a geometric proof of the two-dimensional smooth Poincar\'e
statement. Isothermal coordinates supply the local conformal description,
while uniformization supplies the global identification; the existence of
the local charts alone would not establish the conclusion.

In three dimensions, the analogous statement is that every closed smooth
three-manifold homotopy equivalent to \(S^3\) is diffeomorphic to the
standard three-sphere. Hamilton's approach replaces conformal coordinates
by a geometric heat equation, the Ricci flow,
\[
 \partial_t g=-2\Ric(g).
\]
After fixing the diffeomorphism freedom, this is a nonlinear parabolic
evolution of the metric. Hamilton proved that in dimension three the
volume-normalized flow starting from strictly positive Ricci curvature
converges to a metric of constant positive sectional
curvature~\cite{Hamilton1982}. General initial metrics can develop
singularities, so the broader program required both their analysis and
a way to continue the evolution. Perelman's entropy and noncollapsing
estimates supplied decisive control of singularity models~\cite{Perelman2002}.
His construction of Ricci flow with surgery removes appropriate neck
regions and replaces them by controlled caps~\cite{Perelman2003Surgery}.
For a homotopy three-sphere, finite-time extinction together with the
topology of these surgeries yields the smooth identification with
\(S^3\)~\cite{Perelman2003Extinction}. The lesson for higher dimensions is
that a local smoothing equation must be accompanied by global control of
degeneration and its topological consequences.

In higher dimensions, Smale's handle-cancellation methods, underlying
$h$-cobordism theory, show that every closed smooth homotopy $n$-sphere
with $n\geq5$ is homeomorphic to $S^n$~\cite{Smale1961}.
Milnor's exotic seven-spheres show that the same topological sphere can
admit distinct smooth structures~\cite{Milnor1956}.
Dimension four lies outside the range of these smooth handle-cancellation
arguments. Freedman proved that every closed topological homotopy
four-sphere is homeomorphic to \(S^4\)~\cite{Freedman1982}. The smooth
four-dimensional Poincar\'e conjecture asks whether every closed smooth
homotopy four-sphere is diffeomorphic to the standard \(S^4\). This remains
open. Major approaches include handle calculus and surgery descriptions
that seek to simplify homotopy
spheres, and knot-theoretic obstructions that could distinguish a candidate
exotic sphere~\cite{FreedmanGompfMorrisonWalker2010}.
Geometric approaches seek a distinguished metric or connection
whose existence would force the standard smooth structure. The present
work pursues this direction, with anti-self-dual (ASD) conformal geometry
as a possible target. Fine~\cite{Fine2011} introduced a gradient flow for
definite \(SO(3)\) connections aimed at producing ASD Einstein metrics of
nonzero scalar curvature. This offers a concrete
four-dimensional evolution framework; its application to the smooth
Poincar\'e problem would still require suitable initial connections on
arbitrary homotopy spheres and global control of the flow.

The special geometry of four dimensions explains the appeal of these
variables. Unlike a surface metric, a general four-dimensional metric has
Weyl curvature that cannot be removed by conformal rescaling or a change
of coordinates. At the same time, on an oriented Riemannian four-manifold
the Hodge star satisfies \(*^2=1\) on real two-forms, producing two real
rank-three bundles,
\[
 \Lambda^2T^*M=\Lambda_g^+\oplus\Lambda_g^-.
\]
Atiyah, Hitchin, and Singer~\cite{AHS1978} used this chiral splitting to
connect conformal geometry and gauge fields. In our orientation convention,
their twistor integrability criterion says that \(W^+=0\) precisely when
the natural almost complex
structure on the sphere bundle \(S(\Lambda_g^+)\) is integrable.
The resulting complex threefold encodes a real four-dimensional conformal
geometry through a family of projective lines. Penrose's nonlinear-graviton
construction supplies the corresponding
holomorphic picture and, with additional data, its Ricci-flat
specialization~\cite{Penrose1976}. These are four-dimensional analogues of
the relation between conformal geometry and complex analysis on surfaces,
with a curvature integrability condition in place of an automatic local
complex structure.

Gauge theory provides a complementary reason to work with connections.
Donaldson invariants extract information about smooth four-manifolds from
moduli spaces of anti-self-dual gauge connections~\cite{DonaldsonKronheimer1990}.
Their success motivates the use of gauge fields in smooth topology,
although it does not settle the
homotopy-sphere case, where the second cohomology vanishes. Twistor methods
also reach beyond the self-dual sector. Witten~\cite{Witten1978} describes
the full classical Yang--Mills equations using holomorphic bundles on the
space of complex null lines, with specified triviality and third-order
formal extension conditions. This is a local
complex flat-spacetime correspondence, not an existence theorem for curved
Riemannian four-manifolds. Together these ideas motivate generalized
coordinate descriptions in which maps depend on both a point and a chiral
direction, and curvature is carried by connection data. For a flat-space
precursor based on generalized conformal maps acting on Yang--Mills fields,
see~\cite{Baker2023} and references therein.
Section~\ref{sec:ctx-isothermal} develops this background and its
reconstruction requirements.

For the evolution studied here, the base is a smooth oriented \emph{real}
four-manifold and the metric is positive definite; all perturbative results
concern the closed standard sphere near a round metric. The complex spaces
above supply background and motivation. There are two reasons to develop
variables that remain meaningful away from the integrable twistor setting.
First, an evolution toward an ASD metric must be able to change the Weyl
curvature, rather than assume its desired vanishing at every time. Second,
the round target has nonzero scalar curvature and two nonflat chiral
Levi--Civita connections. Thus setting one entire spin connection to zero,
as can occur in an appropriate Ricci-flat chiral reduction, would exclude
the target itself. Ward's extension of the nonlinear-graviton construction
to Einstein metrics with cosmological constant already distinguishes this
nonzero-scalar-curvature setting from the vacuum one~\cite{Ward1980}.

Definite connections provide such variables on an open set of Riemannian
fields. The Levi--Civita connection induces an \(SO(3)\) connection
\(A_g^+\) on the selfdual bundle \(\Lambda_g^+\).
One may instead start with an independent metric connection \(A\) in an
oriented real rank-three Euclidean bundle \(E\). Using
\(\mathfrak{so}(E)\simeq E\), write its curvature in an oriented
orthonormal internal frame as \(F_A=\sum_iF_i e_i\). We restrict to the
definite region where the matrix \((F_i\wedge F_j)\) is positive relative
to the base orientation. Its curvature span determines a unique
positive-definite conformal metric~\cite{Fine2011}.
With a fixed reference scale \(K>0\), our normalization is
\begin{equation}\label{eq:intro-reconstruction}
 \mu_A=\frac13\sum_iF_i\wedge F_i,\qquad
 \Lambda^+_{\widehat g_A}=\operatorname{span}\{F_1,F_2,F_3\},\qquad
 dV_{\widehat g_A}=\frac{\mu_A}{2K^2}.
\end{equation}
The reconstructed metric \(\widehat g_A\) makes \(F_A\) selfdual.
The further requirement that \(A\) be its chiral Levi--Civita connection
is a differential compatibility condition, not a consequence of definiteness.

Define the symmetric matrix \(Q\) by \(F_i\wedge F_j=Q_{ij}\mu_A\);
then \(Q>0\) and \(\tr Q=3\). Fine's functional, after subtracting its
topological constant, measures the departure from \(Q=I\):
\begin{equation}\label{eq:intro-energy}
 \mathcal P(A)=\int_M\tr((Q-I)^2)\,\mu_A.
\end{equation}
A connection with \(Q=I\) is called \emph{perfect}. In the component
containing the round chiral connection, perfection gives an ASD Einstein
metric with scalar curvature \(12K\)~\cite{Fine2011}. Here ASD means
\(W^+=0\), a statement about the metric's Weyl tensor, distinct from the
selfduality of the gauge curvature \(F_A\).
We study the negative gradient of \(\mathcal P\) for the connection
\(L^2\) pairing formed using \(\widehat g_A\) and the internal metric
\(\delta_{ij}\). It is Fine's flow with a fixed time normalization.
Fine proves short-time existence for definite initial data and, on a compact
four-manifold, convergence modulo gauge and diffeomorphisms for sufficiently
small $L^2_k$ perturbations ($k\geq3$) of a perfect connection~\cite{Fine2011}.
A broader pure-connection formulation of Einstein geometry appears
in~\cite{FineKrasnovPanov2014}.
Our purpose is to make explicit how curvature reconstruction and the
connection flow compare with metric evolution, including the first
nonlinear discrepancies.

Here is the central statement in this setting. Let \(g_0\) be the round
positive-definite metric on the smooth real four-sphere \(S^4\), with
sectional curvature \(K>0\), and fix its orientation. Take
\(E=\Lambda_{g_0}^+\), \(A_0=A_{g_0}^+\), and identify the chiral
bundles of nearby metrics by the positive symmetric coframe identification.
Let \(h\in\Gamma(\operatorname{Sym}^2T^*S^4)\) be a smooth real
transverse-traceless tensor, meaning
\(\nabla^a h_{ab}=0\) and \(\operatorname{tr}_{g_0}h=0\), where
\(\nabla=\nabla^{g_0}\). Infinitesimally, the coframe identification is
\(\dot e^a=\tfrac12h^a{}_b e^b\); it fixes the frame gauge in the
connection variation. For a metric path
\(g_\varepsilon=g_0+\varepsilon h+O(\varepsilon^2)\), define
\[
 a=\mathcal Hh:=\left.\frac{d}{d\varepsilon}\right|_0A_{g_\varepsilon}^+
       \in\Omega^1(S^4,E),\qquad
 \widehat h:=\left.\frac{d}{d\varepsilon}\right|_0
                 \widehat g_{A_{g_\varepsilon}^+}.
\]
Thus \(h\) varies the source metric, \(a\) varies its chiral connection,
and \(\widehat h\) varies the curvature-reconstructed metric; these are
different objects. Put \(L=\nabla^*\nabla\), the nonnegative rough
Laplacian on symmetric two-tensors. The linearized connection flow at
\(A_0\) preserves the image of such tensors under \(\mathcal H\), and
the reconstruction and evolution identities are
\begin{equation}\label{eq:intro-main}
\widehat h=\frac{L+4K}{2K}h,\qquad
\partial_ta=\mathcal H[-4(L+2K)h],\qquad
\partial_t\widehat h=-4(L+2K)\widehat h.
\end{equation}
Here \(t\) is the time of the negative connection-gradient flow just
specified, and the evolution equations are linearizations at \(A_0\),
not equations for arbitrary finite perturbations. The factor four is fixed
by that connection inner product. With \(\tau=4t\), the final equation is
the transverse-traceless linearization of volume-normalized Ricci flow
\(\partial_\tau g=-2\Ric(g)+(\overline s_g/2)g\), where
\(\overline s_g\) is average scalar curvature. The proofs in
Proposition~\ref{rec:linear-reconstruction} and
Section~\ref{sec:ricci-comparison} include the curvature terms as well as
the principal symbols.

The distinction between evolving a connection and evolving a metric is essential. If the same connection energy is first pulled back by \(g\mapsto A_g^+\), its metric gradient has the fourth-order linearization
\begin{equation}\label{eq:intro-fourth}
\partial_th=-(L+2K)(L+4K)h.
\end{equation}
We derive this factorization from the connection calculation and an integration-by-parts identity. It agrees, with the appropriate norm conversion, with the known Hessian of the chiral Weyl energy~\cite{GurskyViaclovsky2015}. The agreement of the two metric energies at quadratic order fails at cubic order through explicit Weyl--Ricci couplings.

The nonlinear reconstruction exhibits a separate issue. Its compatibility with the Levi--Civita connection is measured by the covariant derivative of a curvature-normalized selfdual triple. This quantity vanishes for perfect connections but is generally nonzero away from them. Moreover, the normalized curvature energy is bounded throughout the definite region and does not diverge at its boundary. An exact local family demonstrates how the reconstructed metric can become anisotropic while this energy density remains finite.

The link back to the smooth Poincar\'e question is a conformal rigidity
argument. On a closed smooth homotopy four-sphere, the signature is zero, so an ASD
Riemannian metric has both Weyl components zero. The conformal
developing-map argument then identifies the manifold smoothly with the
standard sphere~\cite{Kuiper1949}. Section~\ref{sec:ctx-topology} gives
the full argument. Obtaining an ASD metric on every such manifold is
therefore an equivalent existence problem. The present calculations describe
the behavior near the round target. They do not construct appropriate
initial definite connections on an arbitrary homotopy four-sphere or
establish the global preservation, compactness and convergence needed to
reach that target. In particular, the transverse-traceless correspondence
does not identify the complete nonlinear connection flow with Ricci flow.

The development is self-contained at the level of the differential-geometric calculations. Standard results used for orientation---the twistor integrability criterion, characteristic-number identities, and the conformal developing-map construction---are stated where needed. We give analytical proofs of the perturbative and reconstruction identities; secondary perturbation examples follow the main evolution theory. No numerical computation is needed for any proposition.

\subsection{Notation and analytical conventions}\label{sec:notation}
All manifolds and fields are smooth and real unless a complexification is explicitly mentioned. The base is an oriented four-manifold, closed whenever an integral is used. The standard sphere is used for perturbation and spectral statements. We use the curvature convention in which a round orthonormal coframe has
\[
R_{ab}=K e^a\wedge e^b,\qquad
\operatorname{Ric}(g_0)=3Kg_0,\qquad s(g_0)=12K.
\]
The scalar Laplacian is \(\Delta=\nabla^a\nabla_a\), while \(L=\nabla^*\nabla=-\Delta\) denotes the nonnegative rough Laplacian on tensors. A tensor \(h\) is transverse-traceless, abbreviated TT, if
\(\nabla^a h_{ab}=0\) and \(\operatorname{tr}_{g_0}h=0\).

The internal bundle \(E\) is oriented, Euclidean, and of rank three, with bracket
\([u,v]_i=\epsilon_{ijk}u_jv_k\). Locally it is the adjoint bundle of a chiral spin bundle. The adjoint formulation itself does not require a global spin structure. We write \(F_A=dA+\tfrac12[A\wedge A]\), or \(dA+A\wedge A\) in matrix notation. Gauge and frame identifications are made explicitly when taking variations.
When spinor notation is used on the real Riemannian base, the spin bundles
have complex rank two and structure group \(SU(2)\). Their connection
forms take values in the real compact Lie algebra \(\mathfrak{su}(2)\);
using complex matrices does not change the field or dimension of the base.
General holomorphic \(SL(2,\mathbb C)\) connections occur only in the
explicitly identified complex background discussion.

Weyl curvature is regarded as an endomorphism of two-forms; its norm is the Hilbert--Schmidt operator norm. The squared norm of the full four-index tensor is four times this squared operator norm. Selfdual triples \(\Sigma_i\) have norm squared \(2\), so
\(\Sigma_i\wedge\Sigma_j=2\delta_{ij}dV\). These normalizations account for the constants in the energy and Hessian identities.

For a symmetric matrix, \([X]_0=X-\tfrac13\operatorname{tr}(X)I\); \(\operatorname{sym}_0\) also takes its symmetric part. The same subscript denotes removal of the metric trace of a symmetric two-tensor. In asymptotic expansions \(X_\varepsilon=\varepsilon X_1+\varepsilon^2X_2+\cdots\), the coefficients do not include factorials. Remainders refer to a smooth family on a fixed compact domain, or to any fixed differentiability norm for which the displayed operations are defined.



\section{Isothermal coordinates and the twistor integrability condition}
\label{sec:ctx-isothermal}

\subsection{Conformal curvature and the Riemannian twistor space}

The surface analogy concerns the choice of conformal variables, rather
than a literal extension of complex coordinates on a Riemann surface.
In real dimension four a conformal change $\widehat g=e^{2u}g$ preserves the
Weyl tensor as a tensor of type $(1,3)$.  It also preserves the Hodge star
on two-forms.  Thus the conditions $W=0$ and $W^+=0$ are properties of
the conformal class, and a flow of the conformal factor cannot turn a
non-ASD class into an ASD class.  Here and below ASD means $W^+=0$;
reversing orientation interchanges the two chiralities.  The corresponding
conformal invariant is
\begin{equation}\label{eq:ctx-Weyl-energy}
 \mathcal E_+(g)=\int_X|W^+_g|^2\,dV_g,
 \qquad \mathcal E_+(e^{2u}g)=\mathcal E_+(g).
\end{equation}
For the signature formula below, the norm is the Hilbert--Schmidt norm
of the Weyl endomorphism of $\Lambda^+$.

The Riemannian twistor space is the sphere bundle
$Z=S(\Lambda_g^+)$, equivalently the bundle of orthogonal complex
structures on the real tangent spaces compatible with the chosen orientation.
The sphere bundle has real dimension six; each fiber is a real two-sphere
with its usual complex structure as $\CP^1$. Its canonical
almost complex structure uses the ordinary complex structure in each
projective-line fiber and the Levi--Civita horizontal distribution.
The Atiyah--Hitchin--Singer integrability theorem states that this almost
complex structure is integrable exactly when $W^+=0$~\cite{AHS1978},
with chirality labels adjusted to our orientation convention. When
integrable, $Z$ is a complex threefold. The base $M$
remains a real four-manifold, and the projection to it is a smooth sphere
bundle projection, not a holomorphic projection to a complex surface.
No Einstein, scalar-curvature, real-analyticity or global spin assumption
is required for this Riemannian integrability criterion.

\subsection{Complex four-dimensional geometry and twistor lines}

For the holomorphic description, work instead on a complex four-manifold
$M_{\mathbb C}$ with a nondegenerate holomorphic symmetric bilinear metric.
Such a metric is complex bilinear, rather than Hermitian; there is no
positive-definite signature in this category. A real-analytic Riemannian
metric has such a description after local complexification. Choose local
complex rank-two spin bundles $S,S'$ with alternating forms and
$TM_{\mathbb C}\simeq S\otimes S'$. For a fixed nonzero primed spinor
$\pi^{A'}$, the vectors $\lambda^A\pi^{A'}$, with $\lambda^A$ arbitrary,
form a complex two-dimensional totally null \emph{alpha plane}.
The correspondence space
\[
 F=\mathbb P(S')\longrightarrow M_{\mathbb C}
\]
therefore has complex dimension five and projective-line fibers. Horizontal
lifts of the alpha planes form a complex rank-two distribution $D_\alpha$.
With the chirality choice used here, its Frobenius integrability is equivalent
to $W^+=0$. Its local leaf space is then a complex threefold~\cite{Penrose1976}.
In the smooth Riemannian description the corresponding involutivity
condition is that of $T^{0,1}Z$ on the real
six-dimensional sphere bundle. These descriptions express the same curvature
criterion in their respective categories; their underlying spaces have
different dimensions and should not be identified.

A point $x$ of complex spacetime is represented in the twistor threefold
by an embedded rational curve $L_x\simeq\CP^1$, whose normal bundle is
\[
 N_{L_x/Z}\simeq\mathcal O(1)\oplus\mathcal O(1).
\]
Here $\mathcal O(k)$ is the degree-$k$ holomorphic line bundle on
$\CP^1$. The identities
$H^1(L_x,N_{L_x/Z})=0$ and $\dim_{\mathbb C}H^0(L_x,N_{L_x/Z})=4$
explain the four complex spacetime parameters: Kodaira's completeness theorem
identifies nearby curves with a local deformation family whose tangent space
is $H^0(N)$~\cite{Kodaira1962}.
His stability theorem explains persistence of the compact curve under
sufficiently small deformations of the ambient complex structure when
$H^1(N)=0$~\cite{Kodaira1963}. Null tangent vectors are the
normal sections that vanish at some point of the curve. This incidence
condition reconstructs the conformal null cone~\cite{Penrose1976}.
Choosing a metric scale and imposing a Riemannian real structure require
further data. The curvature is encoded in
the neighborhood and incidence of the family, not removed by choosing
holomorphic coordinates on a single twistor patch.

\subsection{Direction-dependent maps and connection reconstruction}

This change from points to curves gives a concrete motivation for generalized
isothermal descriptions. In a local twistor chart, write a curve as
$(\zeta,\omega^0(\zeta),\omega^1(\zeta))$ and use flat incidence
$\omega^A=y^{A0'}+\zeta y^{A1'}$. The flat twistor line agreeing with
the curve to first order corresponds to
\[
 y^{A0'}=\omega^A-\zeta\partial_\zeta\omega^A,\qquad
 y^{A1'}=\partial_\zeta\omega^A.
\]
Here the flat conformal metric is represented by
$2(dy^{00'}dy^{11'}-dy^{01'}dy^{10'})$. The formula for $y$ is the affine
tangent-line construction of Shaw~\cite{Shaw1986}, on a patch where the
projective direction can be used as parameter. As the direction varies,
$\partial_\zeta y^{AA'}$ has rank at most one, so it describes a possibly
stationary complex null trajectory. A curved point can thus acquire
different flat representatives in different directions. The natural domain
is correspondence space, and the resulting data must include the direction
and its transition laws. The action on null surfaces and the associated
nonlocal gauge pullback in~\cite{Baker2023} provide the flat-space
motivation for studying such maps.

The connection problem adds a second layer. Ward reconstruction uses a
holomorphic bundle trivial on each \emph{complete} twistor line, together
with the relevant incidence and partial-connection data. The Riemannian
correspondence also requires the compatible positive
reality structure~\cite{AHS1978}. On flat complex correspondence space,
the desired directional operators have the form
\[
 \mathcal D_A(\pi)=\pi^{A'}
       \bigl(\partial_{AA'}+\mathcal A_{AA'}(y)\bigr).
\]
The matrices $\mathcal A_{AA'}(y)$ must depend only on the flat base point.
Regularity on the whole projective direction fiber is what makes global
holomorphic sections of $\mathcal O(1)$ linear in $\pi$; regularity on a
single direction disc would not suffice. This bundle condition is different
from extending one inverse coordinate map over a complete target fiber.
Indeed, if the latter is holomorphic and its source-point projection lies
in one holomorphic coordinate neighborhood, each projected coordinate is
a holomorphic function on $\CP^1$ and hence constant along the fiber.
One must therefore distinguish local geometric charts, globality along the
direction fiber for bundle splitting, and descent to a connection on a
four-dimensional base. An auxiliary bundle obtained by reconstruction
also needs an additional identification with tangent spinors before its
connection can be interpreted as a gravitational spin connection.

The distinction between complex and real geometry is visible in the local
wave expression
\[
 g_H=2\,du\,dv+2\,dz\,dw+H(u,z)\,du^2.
\]
It can be studied holomorphically with complex coordinates and holomorphic
$H$. Its slice with real coordinates and real $H$ has signature $(2,2)$:
both two-by-two metric blocks have negative determinant. Such a slice is
not a positive-definite metric on a real four-manifold. Explicit null
geometry can therefore clarify the coordinate and connection questions
without furnishing the Riemannian initial data required for the sphere
problem. The global exclusion of neutral metrics on a homotopy four-sphere
is proved in Corollary~\ref{cor:ctx-wave-exclusions}.

\subsection{From integrable geometry to a Riemannian flow}

The Einstein condition is another independent issue. The curvature of a
metric's positive chiral Levi--Civita connection contains both a Weyl-plus-scalar
block and a trace-free Ricci block. Selfduality of this connection is
equivalent to the vanishing of the Ricci block, hence to the Einstein
condition~\cite{AHS1978}. It is not equivalent to
$W^+=0$. Penrose's Ricci-flat nonlinear-graviton construction requires
additional twistor structures beyond the conformal line family~\cite{Penrose1976}.
Ward's construction retains Einstein geometry with nonzero cosmological
constant through additional twistor
forms~\cite{Ward1980}. These distinctions matter for $S^4$:
the round metric has $W^+=W^-=0$ and $\Ric=3Kg$, but its two spin
curvatures are $F_i^\pm=-K\Sigma_i^\pm$, as calculated in
Section~\ref{sec:round-model}. Neither chiral connection can be made
zero by a gauge choice when $K>0$.

There is also a classical extension beyond one chiral sector on the gauge
side. The space of complex null lines in flat complex four-space is a
complex five-dimensional \emph{ambitwistor space}. Work on a sufficiently
small complex domain whose intersections with null lines are connected and
simply connected. In the local analytic correspondence, holomorphic bundles
trivial on the $\CP^1\times\CP^1$ of directions
through each point encode gauge connections; the source-free Yang--Mills
equations impose a third-order formal-neighborhood extension
condition~\cite{Witten1978,IsenbergYasskinGreen1978}.
Mason and Skinner~\cite{MasonSkinner2006} review this condition and develop a
formulation on a Cauchy--Riemann submanifold selected by a Euclidean real
slice. The extension and reality conditions are substantive inputs.
These gauge-theoretic constructions do not by
themselves provide a gravitational connection or an evolution of arbitrary
smooth Riemannian metrics.

For the present problem, the immediate obstruction to using an integrable
twistor description throughout a flow is already elementary.

\begin{proposition}[An obstruction to exact flat distribution charts]
\label{prop:ctx-flat-distribution}
Suppose a complexified correspondence-space chart is a local
biholomorphism carrying $D_\alpha$ to the alpha distribution of flat
space.  Then $W^+=0$ on the projected source domain.  In the smooth
Riemannian formulation, an exact local identification of the canonical
twistor almost complex structure with an integrable one has the same
consequence.
\end{proposition}
\begin{proof}
For any distribution $D$, its Frobenius tensor is
\[
 \mathfrak F_D:\Lambda^2D\longrightarrow T/D,
 \qquad \mathfrak F_D(U,V)=[U,V]\bmod D.
\]
Lie brackets commute with pushforward by a local diffeomorphism.
Consequently an exact equivalence of distributions preserves the
vanishing of $\mathfrak F_D$.  The flat alpha distribution is
involutive, so the source one is involutive.  The curvature integrability
criterion gives $W^+=0$.  The smooth assertion follows in the same way
from naturality of the Nijenhuis tensor.
\end{proof}

An enlarged system intended to evolve arbitrary Riemannian metrics
must therefore retain the nonintegrability tensor, or another variable
encoding $W^+$.  It cannot be an exact flat-distribution equivalence
at every time.  In particular, a construction that starts with an
integrable curved twistor space describes an already restricted
conformal class; its holomorphic coordinates do not by themselves
provide coordinates for a general metric flow.



\section{Chiral Riemannian geometry and the round model}
\label{sec:round-model}

The round four-sphere supplies a nonzero-scalar-curvature model in which
connection descent, metric reconstruction, and global bundle patching can all
be examined explicitly.  Its role is structural: it identifies which features
of a chiral description persist in positive-definite geometry and which
features of a vacuum reduction must be abandoned.

\subsection{Curvature conventions and the two chiral connections}
\label{subsec:round-conventions}

Let $(M,g)$ be an oriented Riemannian spin four-manifold; all formulas in this
subsection are also valid locally without a global spin assumption.  Choose
an oriented orthonormal coframe $(e^1,e^2,e^3,e^4)$, with
$e^{1234}=dV_g>0$.  The Levi--Civita connection and curvature forms satisfy
\begin{equation}
 de^a+\omega^a{}_b\wedge e^b=0,
 \qquad R^a{}_b=d\omega^a{}_b+\omega^a{}_c\wedge\omega^c{}_b.
 \label{eq:round-cartan-conventions}
\end{equation}
Our sectional-curvature convention is
$R_{ijkl}=\langle R(e_i,e_j)e_l,e_k\rangle$, so a metric of constant
sectional curvature $K$ has $R_{1212}=K$ and
$R_{ab}=K e^a\wedge e^b$.

For $i,j,k\in\{1,2,3\}$, define the two sets of symbols
\begin{equation}
 \eta^\pm_{i,jk}=\epsilon_{ijk},\qquad
 \eta^\pm_{i,j4}=\pm\delta_{ij},\qquad
 \eta^\pm_{i,4j}=\mp\delta_{ij},\qquad
 \eta^\pm_{i,44}=0.
 \label{eq:round-eta}
\end{equation}
Set
\begin{equation}
 \Sigma_i^\pm=\frac12\eta^\pm_{i,ab}e^a\wedge e^b,
 \qquad T_i=-\frac{i}{2}\sigma_i,
 \label{eq:round-sigma-generators}
\end{equation}
where $\sigma_i$ are the Pauli matrices.  Then
\begin{align}
 *\Sigma_i^\pm&=\pm\Sigma_i^\pm,
 &\langle\Sigma_i^\pm,\Sigma_j^\pm\rangle&=2\delta_{ij},
 \label{eq:round-sigma-norms}\\
 \Sigma_i^\pm\wedge\Sigma_j^\pm&=\pm2\delta_{ij}dV_g,
 &\Sigma_i^+\wedge\Sigma_j^-&=0,
 \label{eq:round-sigma-wedges}\\
 [T_i,T_j]&=\epsilon_{ijk}T_k,
 &\operatorname{tr}(T_iT_j)&=-\frac12\delta_{ij}.
 \label{eq:round-generator-norms}
\end{align}
As skew $4\times4$ matrices, the $\eta_i^\pm$ obey
$[\eta_i^\pm,\eta_j^\pm]=-2\epsilon_{ijk}\eta_k^\pm$, and the two
sets commute with one another.  Thus $-\eta_i^\pm/2$ give the two
$\mathfrak{su}(2)$ ideals in $\mathfrak{so}(4)$.  With these conventions,
the two spin connections are
\begin{equation}
 A^\pm=A_i^\pm T_i,\qquad
 A_i^\pm=-\frac12\eta^\pm_{i,ab}\omega_{ab},\qquad
 F_i^\pm=-\frac12\eta^\pm_{i,ab}R_{ab}.
 \label{eq:round-spin-projection}
\end{equation}
The last equality follows because the projections are Lie algebra
homomorphisms, with $F^\pm=dA^\pm+A^\pm\wedge A^\pm$.

Use $\Sigma_i^\pm/\sqrt2$ as orthonormal bases of $\Lambda^\pm$.
The curvature operator has the block decomposition
\begin{equation}
 \mathcal R_g=
 \begin{pmatrix}
 C_+&B\\ B^T&C_-
 \end{pmatrix},
 \qquad C_\pm=W^\pm+\frac{s}{12}I.
 \label{eq:round-curvature-blocks}
\end{equation}
Here $s$ is scalar curvature and $B:\Lambda^-\to\Lambda^+$ is the
block corresponding to trace-free Ricci curvature.  In particular, $B=0$
if and only if $g$ is Einstein.  Equation~\eqref{eq:round-spin-projection}
becomes
\begin{align}
 F_i^+&=-(C_+)_{ij}\Sigma_j^+-B_{ij}\Sigma_j^-,
 \label{eq:round-plus-block}\\
 F_i^-&=-(B^T)_{ij}\Sigma_j^+-(C_-)_{ij}\Sigma_j^-.
 \label{eq:round-minus-block}
\end{align}
The labels $+$ and $-$ refer to form chirality.  Thus $A^+$ is self-dual
as a gauge connection precisely when $B=0$, whereas $A^-$ is anti-self-dual
precisely when $B=0$.  This Einstein condition is distinct from the
half-conformal-flatness condition $W^+=0$.  All norms of $W^\pm$ used here
are Hilbert--Schmidt norms of curvature operators: the squared four-index
tensor norm is four times the squared operator norm.

\subsection{An explicit positive-curvature model}
\label{subsec:round-explicit}

Let $g_0$ be the round metric of radius $\ell>0$ on $S^4$, and write
$K=\ell^{-2}$.  On a stereographic patch, put
\begin{equation}
 r^2=\sum_{a=1}^4(x^a)^2,\qquad
 D=\ell^2+r^2,\qquad
 \Omega=\frac{2\ell^2}{D},\qquad
 g_0=\Omega^2\sum_{a=1}^4(dx^a)^2,
 \qquad e^a=\Omega\,dx^a.
 \label{eq:round-stereographic}
\end{equation}

\begin{proposition}[Round chiral connections]
\label{prop:round-chiral-connections}
In the coframe~\eqref{eq:round-stereographic},
\begin{align}
 \omega_{ab}&=\frac{2(x_a\,dx_b-x_b\,dx_a)}{D},
 &R_{ab}&=K e^a\wedge e^b,
 \label{eq:round-omega-curvature}\\
 A_i^\pm&=\frac{2\eta^\pm_{i,ab}x_b\,dx^a}{D},
 &F_i^\pm&=-K\Sigma_i^\pm.
 \label{eq:round-A-F}
\end{align}
Consequently $\operatorname{Ric}(g_0)=3K g_0$, $s(g_0)=12K$, and
$W^+=W^-=0$, while both spin connections have nonzero curvature.  Relative
to the flat stereographic metric as well as to $g_0$, their curvatures obey
$*F^\pm=\pm F^\pm$.
\end{proposition}

\begin{proof}
Since $d\log\Omega=-2x_a dx^a/D$, the displayed skew connection matrix
satisfies the torsion equation~\eqref{eq:round-cartan-conventions}.  It is
therefore the Levi--Civita connection.  Write $\theta=x_a dx^a$ and
$\alpha_{ab}=x_a dx_b-x_b dx_a$.  Direct differentiation and multiplication
give
\begin{align*}
 d\omega_{ab}
 &=\frac4D dx_a\wedge dx_b-
   \frac4{D^2}\theta\wedge\alpha_{ab},\\
 \omega_{ac}\wedge\omega_{cb}
 &=\frac4{D^2}\bigl(
   \theta\wedge\alpha_{ab}-r^2dx_a\wedge dx_b\bigr).
\end{align*}
Their sum is
\[
 R_{ab}=\frac{4\ell^2}{D^2}dx_a\wedge dx_b
       =\ell^{-2}e^a\wedge e^b.
\]
Projecting by~\eqref{eq:round-spin-projection} proves both chiral formulas.
In coordinate components they read
\begin{equation}
 F_{i,\mu\nu}^\pm=-\frac{4\ell^2}{D^2}\eta^\pm_{i,\mu\nu},
 \qquad
 F_i^\pm=\frac12F_{i,\mu\nu}^\pm dx^\mu\wedge dx^\nu.
 \label{eq:round-F-components}
\end{equation}
The curvature identities give the stated Ricci, scalar, and Weyl tensors.
Finally, Hodge star on two-forms in four dimensions is unchanged by conformal
rescaling, so the chiral gauge equations hold for the flat metric as well.
\end{proof}

The quadratic term \(A^\pm\wedge A^\pm\) is essential: these are nonabelian chiral fields. The stereographic chart represents them as ordinary selfdual and anti-selfdual gauge connections on a flat coordinate domain. It does not make either connection flat. Indeed each \(\Sigma_i^\pm\) is a nondegenerate two-form, so the common curvature kernel is zero. A description that imposes flatness of the opposite chiral connection would therefore exclude the round target.

\subsection{Global patching and characteristic forms}
\label{subsec:round-global}

Identify $x$ with the quaternion
$x=x^4+x^1\mathbf i+x^2\mathbf j+x^3\mathbf k$, and identify the imaginary
quaternion units with $2T_1,2T_2,2T_3$.  The local formulas take the form
\begin{equation}
 A^+=\frac{\operatorname{Im}(\bar x\,dx)}{\ell^2+|x|^2},
 \qquad
 A^-=\frac{\operatorname{Im}(x\,d\bar x)}{\ell^2+|x|^2}.
 \label{eq:round-quaternion-connections}
\end{equation}
The second chart is $y=\ell^2x^{-1}$.  This is an orientation-preserving
transition between oriented stereographic coordinate domains, and the round
metric has the same expression~\eqref{eq:round-stereographic} in $y$.

\begin{proposition}[Extension across the omitted pole]
\label{prop:round-patching}
On the overlap, put $q=x/|x|$.  With gauge convention
$A^h=h^{-1}Ah+h^{-1}dh$, the changes
\begin{equation}
 h^+=q^{-1},\qquad h^-=q
 \label{eq:round-clutching}
\end{equation}
transform~\eqref{eq:round-quaternion-connections} into the same regular
formulas in $y$.  Thus both fields extend smoothly as connections on their
chiral spin bundles over $S^4$.
\end{proposition}

\begin{proof}
Write $x=rq$, and set $f=r^2/(\ell^2+r^2)$.  Then
$A^+=f q^{-1}dq$ and $A^-=f q\,d(q^{-1})$.  Their gauge transforms are
\[
 (A^+)^{q^{-1}}=-(1-f)dq\,q^{-1},
 \qquad
 (A^-)^q=(1-f)q^{-1}dq.
\]
On the other hand $y=(\ell^2/r)q^{-1}$, and
\[
 \frac{|y|^2}{\ell^2+|y|^2}
 =\frac{\ell^2}{\ell^2+r^2}=1-f.
\]
Substitution into the two expressions in~\eqref{eq:round-quaternion-connections}
gives exactly the transformed fields.  They are smooth at $y=0$.
These gauge factors are the spin lifts, up to constant central signs, of
the coframe transition: in quaternion notation it is
$e_y=-q^{-1}e_xq^{-1}$.  Thus the extension agrees with the global
Levi--Civita spin connections.
\end{proof}

Use the Chern--Weil convention
\begin{equation}
 c_2(S^\pm)[S^4]
 =\frac1{8\pi^2}\int_{S^4}\operatorname{tr}(F^\pm\wedge F^\pm).
 \label{eq:round-c2-convention}
\end{equation}
Equations~\eqref{eq:round-sigma-wedges} and
\eqref{eq:round-generator-norms} give
\begin{equation}
 \operatorname{tr}(F^\pm\wedge F^\pm)=\mp3K^2dV_{g_0},
 \qquad
 -\operatorname{tr}(F^\pm\wedge *F^\pm)=3K^2dV_{g_0}.
 \label{eq:round-characteristic-densities}
\end{equation}
Since $\operatorname{Vol}(S^4,g_0)=8\pi^2/(3K^2)$, the chosen labels give
\begin{equation}
 c_2(S^+)[S^4]=-1,\qquad c_2(S^-)[S^4]=+1,
 \qquad
 -\int_{S^4}\operatorname{tr}(F^\pm\wedge *F^\pm)=8\pi^2.
 \label{eq:round-charge-energy}
\end{equation}
The signs of the characteristic numbers depend on the stated orientation
and labels; the nontriviality of both bundles does not.  A regular local
connection form therefore cannot replace the nontrivial overlap data.

For later comparison with definite connections, the positive four-form
associated to $A^+$ in our component normalization is
\begin{equation}
 \mu_{A^+}:=\frac13\sum_iF_i^+\wedge F_i^+
 =2K^2dV_{g_0},
 \qquad
 \int_{S^4}\mu_{A^+}=\frac{16\pi^2}{3}.
 \label{eq:round-curvature-volume}
\end{equation}
It is independent of the radius after integration.  Using the negative of
the corresponding expression for $A^-$ gives the same positive four-form.
The span of the three $F_i^+$ recovers $\Lambda^+$ and hence the conformal
class of $g_0$; a volume convention fixes a representative.  In particular,
the curvature-volume convention~\eqref{eq:round-curvature-volume} fixes a
normalization, rather than recovering an independently specified original
radius.

\subsection{Scale normalization}
\label{subsec:round-scale-twistor}

There is a separate loss of scale in a connection-only description.
Multiplying a metric by a positive spatially constant factor leaves its
Levi--Civita connection unchanged.  For the round sphere, unnormalized
Ricci flow is
\begin{equation}
 g(t)=(1-6Kt)g_0,\qquad 0\leq t<(6K)^{-1}.
 \label{eq:round-ricci-flow}
\end{equation}
Its two spin connections remain fixed under the natural identifications
of the orthonormal frame bundles, even though the sphere shrinks.  Any
comparison between metric evolution and connection evolution must
therefore supply a scale variable or a scale normalization.  This issue
is already present in the exact round solution and is independent of
perturbative or singularity questions.


\section{Curvature reconstruction and its differential}
\label{rec:section}

Let $M$ be an oriented smooth four-manifold and let $E\to M$ be an
oriented Euclidean bundle of rank three.  Its oriented orthonormal frames
identify its Lie bracket with
$[e_i,e_j]=\epsilon_{ijk}e_k$.  A metric connection $A$ on $E$ has curvature
$F_A=F_i e_i\in\Omega^2(M,E)$.  Repeated internal indices are summed.
The normalization of the internal metric is the one for which
$\langle e_i,e_j\rangle=\delta_{ij}$; in an $\mathfrak{su}(2)$ realization
it is $-2\operatorname{tr}(T_iT_j)=\delta_{ij}$.

We call $A$ \emph{definite with positive curvature plane} if the matrix of four-forms
$(F_i\wedge F_j)$ is positive definite relative to the prescribed
orientation.  Equivalently, for any positive four-form $\nu$, the matrix
$(F_i\wedge F_j)/\nu$ is positive definite.  Define
\begin{equation}
 \mu_A=\frac13\sum_iF_i\wedge F_i,
 \qquad F_i\wedge F_j=Q_{ij}\mu_A.
 \label{rec:mu-Q}
\end{equation}
Then $\mu_A>0$, $Q=Q^T>0$, and $\operatorname{tr}Q=3$.  The four-form is frame independent and the matrix is equivariant under internal orthonormal changes of frame. This wedge-positivity convention is distinct from the orientation sign of the curvature map, which is fixed separately below.

\subsection{Positive three-planes and the graph map}

We first record the linear algebra, including the normalization that will
enter all variation formulas.  Fix an oriented Euclidean four-space
$(V,g)$ and an oriented orthonormal coframe $(e^1,e^2,e^3,e^4)$, and use
the two-form bases
\begin{equation}
 \begin{split}
 (\Sigma_1^+,\Sigma_2^+,\Sigma_3^+)
   &=(e^{12}+e^{34},\ e^{13}-e^{24},\ e^{14}+e^{23}),\\
 (\Sigma_1^-,\Sigma_2^-,\Sigma_3^-)
   &=(e^{12}-e^{34},\ e^{13}+e^{24},\ e^{14}-e^{23}).
 \end{split}
 \label{rec:form-bases}
\end{equation}
Thus $|\Sigma_i^\pm|^2=2$ and
\begin{equation}
 \Sigma_i^\pm\wedge\Sigma_j^\pm
    =\pm2\delta_{ij}\,dV_g,
 \qquad \Sigma_i^+\wedge\Sigma_j^-=0.
 \label{rec:wedge-normalization}
\end{equation}
In particular, $\Sigma_i^\pm/\sqrt2$ are orthonormal; the unnormalized
bases in \eqref{rec:form-bases} are used in all matrix formulas below.
We use a row convention for graph matrices: a matrix $T=(T_{ij})$
represents the plane spanned by
$\Sigma_i^++T_{ij}\Sigma_j^-$.  It represents a linear map from
$\Lambda_g^+$ to $\Lambda_g^-$, with the source index written first.

For a $g$-self-adjoint trace-free endomorphism $H$ of $V$, set
$g_H=g(e^H\,\cdot,\cdot)$.  The positive symmetric coframe
identification is $e_H=e^{H/2}e$.  Its infinitesimal action on two-forms is
the self-adjoint operator
\begin{equation}
 \mathscr D_H(\alpha\wedge\gamma)
   =\frac12\bigl(H^*\alpha\wedge\gamma
                         +\alpha\wedge H^*\gamma\bigr).
 \label{rec:D-H}
\end{equation}
Define $\beta_g(H)$ by
\begin{equation}
 \mathscr D_H\Sigma_i^+=\beta_g(H)_{ij}\Sigma_j^-.
 \label{rec:beta-definition}
\end{equation}
The absence of a self-dual component in this definition is part of the
following lemma.  When applying $\beta_g$ to a trace-free symmetric
two-tensor, we raise one index using $g$.

\begin{lemma}[The logarithmic graph map]
\label{rec:graph-lemma}
The map
\[
 \beta_g:\operatorname{Sym}_0(V,g)\longrightarrow
                  \operatorname{Hom}(\Lambda_g^+,\Lambda_g^-)
\]
is an isomorphism and satisfies
\begin{equation}
 |\beta_g(H)|^2=\frac14|H|^2.
 \label{rec:beta-norm}
\end{equation}
For diagonal $H=\operatorname{diag}(h_1,h_2,h_3,h_4)$ with
$\sum_a h_a=0$, the bases \eqref{rec:form-bases} give
\begin{equation}
 \beta_g(H)=\operatorname{diag}
 \left(\frac{h_1+h_2}{2},\frac{h_1+h_3}{2},
                         \frac{h_1+h_4}{2}\right).
 \label{rec:beta-diagonal}
\end{equation}
If $T(H)$ is the graph matrix of $\Lambda_{g_H}^+$ over
$\Lambda_g^+$, then
\begin{equation}
 T(H)=\beta_g(H)-\frac13\beta_g(H)\beta_g(H)^T\beta_g(H)
                         +O(|H|^5).
 \label{rec:graph-cubic}
\end{equation}
In particular, the quadratic term of the logarithmic graph map vanishes.
\end{lemma}

\begin{proof}
Diagonalize $H$ in an oriented orthonormal coframe.  The operator
$\mathscr D_H$ multiplies $e^{ab}$ by $(h_a+h_b)/2$.  Complementary pairs
have opposite eigenvalues because $\operatorname{tr}H=0$.  It follows
directly that $\mathscr D_H$ interchanges $\Lambda_g^+$ and
$\Lambda_g^-$ and that its matrix, acting on the ordered column of basis
forms, is
\begin{equation}
 \mathscr D_H\begin{pmatrix}\Sigma^+\\\Sigma^-\end{pmatrix}
 =\begin{pmatrix}0&\beta_g(H)\\\beta_g(H)^T&0\end{pmatrix}
                       \begin{pmatrix}\Sigma^+\\\Sigma^-\end{pmatrix}.
 \label{rec:block-generator}
\end{equation}
The lower-left block is the transpose because $\mathscr D_H$ is
self-adjoint and both bases have the same norm.  Formula
\eqref{rec:beta-diagonal} follows.  In that formula
$\sum_{j=2}^4(h_1+h_j)^2/4=\sum_{a=1}^4h_a^2/4$.
Orthogonal equivariance proves \eqref{rec:beta-norm} for every $H$.
Both spaces have dimension nine, so the norm identity proves that
$\beta_g$ is an isomorphism.

The full coframe action on two-forms is $\exp(\mathscr D_H)$.  Writing
$b=\beta_g(H)$, its first three rows have the form
$A\Sigma^++B\Sigma^-$, where
\[
 A=I+\frac12bb^T+O(|b|^4),\qquad
 B=b+\frac16bb^Tb+O(|b|^5).
\]
Consequently the graph matrix is $A^{-1}B$, which gives
\eqref{rec:graph-cubic}.  More exactly,
\[
 T(H)=
 \frac{\tanh\sqrt{bb^T}}{\sqrt{bb^T}}\,b,
\]
where the quotient has its analytic value $1$ at zero.  This also
exhibits the oddness of the graph map in $H$.
\end{proof}

\begin{lemma}[Conformal metric from a positive plane]
\label{rec:positive-plane}
A positive three-plane in $\Lambda^2V^*$, for the wedge pairing and a
fixed orientation of $V$, is $\Lambda_{\widetilde g}^+$ for a unique
positive-definite conformal class $[\widetilde g]$.
\end{lemma}

\begin{proof}
Choose a reference metric $g$.  Projection of a positive plane $P$ to
$\Lambda_g^+$ is injective: a nonzero element in its kernel would lie in
the negative plane $\Lambda_g^-$.  The projection is therefore an
isomorphism.  Write $P$ as the graph of $T$.  By
\eqref{rec:wedge-normalization}, positivity is equivalent to
$I-TT^T>0$, or to all singular values of $T$ being less than one.

Independent rotations of the two chiral bases can be induced by an
oriented orthogonal change of coframe.  Indeed the action
$\operatorname{SO}(4)\to\operatorname{SO}(3)\times\operatorname{SO}(3)$
has differential given by the two chiral summands of
$\mathfrak{so}(4)$; that differential is an isomorphism.  Its image is
thus open, and it is all of the connected target.  Apply a signed
singular-value decomposition to arrange
$T=\operatorname{diag}(t_1,t_2,t_3)$, with $|t_i|<1$, in such coframes.
Put $r_i=\operatorname{arctanh}t_i$ and set
\[
 (h_1,h_2,h_3,h_4)
  =(r_1+r_2+r_3,\ r_1-r_2-r_3,
                  -r_1+r_2-r_3,\ -r_1-r_2+r_3).
\]
Then $\operatorname{tr}H=0$ and \eqref{rec:beta-diagonal} gives
$\beta_g(H)=\operatorname{diag}(r_1,r_2,r_3)$.
The exact graph formula in Lemma~\ref{rec:graph-lemma} gives
$T(H)=T$, establishing existence with $\widetilde g=g(e^H\cdot,\cdot)$.
Equivalently, without choosing singular vectors, the inverse formula is
\[
 \beta_g(H)=
   \frac{\operatorname{arctanh}\sqrt{TT^T}}{\sqrt{TT^T}}\,T.
\]
The quotient is analytic at zero and on the domain $I-TT^T>0$.
This also proves smooth dependence of the normalized metric on $P$.

For uniqueness, suppose two metrics have the same self-dual plane.
Their anti-self-dual planes are its wedge-orthogonal complement, so
their Hodge stars on two-forms agree.  Diagonalize one metric relative
to the other, with positive eigenvalues $\lambda_1,\ldots,\lambda_4$.
Equality of the self-dual planes implies
\[
 \lambda_1\lambda_2=\lambda_3\lambda_4,\quad
 \lambda_1\lambda_3=\lambda_2\lambda_4,\quad
 \lambda_1\lambda_4=\lambda_2\lambda_3.
\]
Positivity implies $\lambda_1=\lambda_2=\lambda_3=\lambda_4$.
The metrics are conformally related.
\end{proof}

\subsection{Reconstruction with a fixed scale}

Fix a constant $K>0$.  Apply Lemma~\ref{rec:positive-plane} to the span
of $F_1,F_2,F_3$, and fix its conformal scale by
\begin{equation}
 \Lambda_{\widehat g_A}^+=\operatorname{span}\{F_1,F_2,F_3\},
 \qquad dV_{\widehat g_A}=\frac{\mu_A}{2K^2}.
 \label{rec:metric-definition}
\end{equation}
Uniqueness makes this construction global.  It is smooth on the open
set of definite connections, as follows from the graph
description and the positive matrix functions used above.  If $g_0$ is
round with sectional curvature $K$ and its chiral connection satisfies
$F_{0i}=-K\Sigma_{0i}^+$, then
$\mu_{A_0}=2K^2dV_{g_0}$ and $\widehat g_{A_0}=g_0$.

There are two choices for the orientation of the curvature map
$E\to\Lambda_{\widehat g_A}^+$.  From now on we use the component
containing the round connection just described.  Its curvature map
reverses the internal orientation.  Define
\begin{equation}
 S=Q^{1/2},\qquad
 \Sigma_A=-K^{-1}S^{-1}F_A.
 \label{rec:adapted-frame}
\end{equation}
Then
\begin{equation}
 F_A=-KS\Sigma_A,
 \qquad (\Sigma_A)_i\wedge(\Sigma_A)_j
       =2\delta_{ij}\,dV_{\widehat g_A}.
 \label{rec:adapted-frame-wedge}
\end{equation}
Thus $\Sigma_A/\sqrt2$ identifies $E$ isometrically and with the
appropriate orientation with $\Lambda_{\widehat g_A}^+$.  This is an
adapted bundle identification, not an assertion that a prescribed
self-dual coframe can always be put into this form by an internal
gauge transformation alone.  All formulas are local in frames and
patch tensorially.

\subsection{Exact metric velocity}

\begin{proposition}[Metric and normalized-curvature velocities]
\label{rec:velocity-proposition}
Let $A(t)$ be a smooth path of definite connections in the
chosen component, with $v=\partial_tA$.  In the adapted self-dual frame
$\Sigma_A$ and any normalized anti-self-dual frame $\Sigma^-$ for
$\widehat g_A$, write
\begin{equation}
 d_Av=U\Sigma_A+V\Sigma^-.
 \label{rec:velocity-decomposition}
\end{equation}
Let $h=\partial_t\widehat g_A$ and
$h_0=\widehat g_A^{-1}h-\tfrac14\operatorname{tr}_{\widehat g_A}(h)I$.
Then
\begin{align}
 \beta_{\widehat g_A}(h_0)&=-K^{-1}S^{-1}V,
 \label{rec:velocity-traceless}\\
 \operatorname{tr}_{\widehat g_A}h
       &=-\frac4{3K}\operatorname{tr}(SU),
 \label{rec:velocity-trace}\\
 \partial_tQ
       &=-K^{-1}(US+SU^T)
           +\frac2{3K}\operatorname{tr}(SU)Q.
 \label{rec:velocity-Q}
\end{align}
In particular,
\begin{equation}
 \partial_t\widehat g_A
 =\widehat g_A\left(
       -K^{-1}\beta_{\widehat g_A}^{-1}(S^{-1}V)
       -\frac1{3K}\operatorname{tr}(SU)I
                  \,\cdot,\cdot\right).
 \label{rec:velocity-metric}
\end{equation}
Equation~\eqref{rec:velocity-Q} uses a fixed internal frame; in a
time-dependent internal frame its left side is the corresponding
covariant time derivative.
\end{proposition}

\begin{proof}
Differentiate $F=-KS\Sigma_A$.  Changing the normalization or rotating
the self-dual frame produces only self-dual terms.  By the definition
of $\beta$, the anti-self-dual component of
$\partial_t\Sigma_A$ is $\beta(h_0)\Sigma^-$.  Since
$\partial_tF=d_Av$, comparison of the anti-self-dual components gives
$V=-KS\beta(h_0)$, proving \eqref{rec:velocity-traceless}.

For the remaining identities, differentiation of
$F_i\wedge F_j=Q_{ij}\mu_A$ and
\eqref{rec:adapted-frame-wedge} gives
\[
 (\partial_tQ)\mu_A+Q\partial_t\mu_A
       =-2K(US+SU^T)dV_{\widehat g_A}
       =-K^{-1}(US+SU^T)\mu_A.
\]
Taking the trace and using $\operatorname{tr}Q=3$ yields
\begin{equation}
 \frac{\partial_t\mu_A}{\mu_A}
       =-\frac2{3K}\operatorname{tr}(SU).
 \label{rec:mu-velocity}
\end{equation}
Because $dV_{\widehat g_A}=\mu_A/(2K^2)$ and $K$ is fixed, the
left side is $\tfrac12\operatorname{tr}_{\widehat g_A}h$.
This proves \eqref{rec:velocity-trace}, and substitution back into
the matrix identity proves \eqref{rec:velocity-Q}.
\end{proof}

For an invariant version of the same calculation, set
$Z_{ij}=((d_Av)_i\wedge F_j)/\mu_A$.  Then $Z=-US/K$ and
\begin{equation}
 \partial_tQ=Z+Z^T-\frac23\operatorname{tr}(Z)Q,
 \qquad
 \partial_t\mu_A=\frac23\operatorname{tr}(Z)\mu_A
                =\frac23d\left(\sum_i v_i\wedge F_i\right).
 \label{rec:invariant-velocities}
\end{equation}
The last equality follows from $d_AF_A=0$.
If $M$ is closed, the reconstructed total volume is therefore constant
along every such path on the fixed bundle.  This fixes total scale,
not the pointwise volume density.  For the round component on $S^4$,
\[
 \int_{S^4}\mu_A=\frac{16\pi^2}{3},\qquad
 \operatorname{Vol}(S^4,\widehat g_A)=\frac{8\pi^2}{3K^2}.
\]
The anti-self-dual matrix $V$ controls the trace-free metric velocity
but does not occur directly in $\partial_tQ$.  Thus these identities
alone do not turn a bound on $Q-I$ into a bound on metric derivatives.

\subsection{The Levi--Civita compatibility defect}

Use \eqref{rec:adapted-frame} to transport the chiral Levi--Civita
connection of $\widehat g_A$ to $E$, and denote the resulting connection
by $\Gamma_A^+$.  In particular, $d_{\Gamma_A^+}\Sigma_A=0$.

\begin{proposition}[Compatibility with the reconstructed metric]
\label{rec:compatibility}
The $E$-valued three-form
\begin{equation}
 \mathcal T_A=d_A\Sigma_A
   =-K^{-1}(d_AQ^{-1/2})\wedge F_A
 \label{rec:torsion-defect}
\end{equation}
determines the difference $b=A-\Gamma_A^+$.  More precisely,
\begin{equation}
 \mathcal T_A=[b,\Sigma_A],
 \qquad \mathcal T_A=0\ \Longleftrightarrow\ A=\Gamma_A^+.
 \label{rec:compatibility-equivalence}
\end{equation}
With the stated form and internal norms, one also has the pointwise
estimate
\begin{equation}
 |b|\leq |\mathcal T_A|\leq2|b|.
 \label{rec:compatibility-norm}
\end{equation}
\end{proposition}

\begin{proof}
Differentiate \eqref{rec:adapted-frame} covariantly and use the Bianchi
identity to obtain \eqref{rec:torsion-defect}.  Subtracting
$d_{\Gamma_A^+}\Sigma_A=0$ gives the first identity in
\eqref{rec:compatibility-equivalence}.

For completeness, the needed pointwise invertibility can be checked
without a differential equation.  Choose a coframe in which
$(\Sigma_A)_{i,ab}=\eta_{i,ab}$, with
\[
 \eta_{i,jk}=\epsilon_{ijk},\qquad
 \eta_{i,j4}=\delta_{ij},\qquad
 \eta_{i,4j}=-\delta_{ij}.
\]
Define an endomorphism of $E$-valued one-forms by
$\mathscr Cb=*_{\widehat g_A}[b,\Sigma_A]$.  Its components are
\[
 (\mathscr Cb)_{ic}
       =\epsilon_{ijk}\eta_{k,dc}b_{jd}.
\]
The elementary contraction identity
\[
 \eta_{i,ab}\eta_{j,cb}
       =\delta_{ij}\delta_{ac}+\epsilon_{ijk}\eta_{k,ac}
\]
gives $\mathscr C^2+\mathscr C-2I=0$.  Also $\mathscr C$ is
self-adjoint, since both the internal epsilon tensor and the eta
matrix change sign when their paired indices are exchanged.
Its only possible eigenvalues are $1$ and $-2$, so it is invertible and
\[
 \mathscr C^{-1}=\frac12(\mathscr C+I),\qquad
 b=\frac12(\mathscr C+I)*_{\widehat g_A}\mathcal T_A.
\]
This proves the equivalence and the norm bounds.
\end{proof}

If $Q=I$, then $\mathcal T_A=0$, and the reconstructed connection has
$F_A=-K\Sigma_A$.  The usual chiral curvature decomposition then gives
$\Ric_0(\widehat g_A)=0$, $W^+(\widehat g_A)=0$, and
$s(\widehat g_A)=12K$.  More generally, $\mathcal T_A=0$ implies that
$\widehat g_A$ is Einstein: the curvature of its chiral Levi--Civita
connection is self-dual by construction, so its opposite-chirality
Ricci block vanishes.

This compatibility condition is stronger than requiring that $A$ be
the chiral Levi--Civita connection of some other metric $g$.
It does not assert that an arbitrary $A_g^+$ is Levi--Civita for its
curvature-reconstructed metric.  It instead gives an exact additional
variable that measures this discrepancy.  In particular, the formulas
above do not establish a closed nonlinear evolution equation in
$\widehat g_A$ alone.

\section{Reconstruction from a metric's chiral connection}
\label{rec:source-metric-section}

We now start with a Riemannian metric $g$ on $S^4$ near its round metric
$g_0$ of sectional curvature $K$.  Use its oriented orthonormal coframes to
identify the internal bundle with its chiral bundle, and write
\begin{equation}
 F_{A_g^+}=-C\Sigma_g^+-B\Sigma_g^-,\qquad
 C=W^++\frac{s}{12}I.
 \label{rec:chiral-curvature}
\end{equation}
Here $W^+$ is the curvature-operator Weyl block, so its norm is the
Hilbert--Schmidt norm of a three-by-three symmetric trace-free matrix.
For either Weyl component, the squared four-index tensor norm is four
times the squared curvature-operator norm.  The sign in
\eqref{rec:chiral-curvature} agrees with $F_{A_0}=-K\Sigma_0^+$.

\begin{lemma}[The Ricci block and the curvature graph]
\label{rec:ricci-block}
With the row and form conventions above,
\begin{equation}
 B=\beta_g(g^{-1}\Ric_0(g)),\qquad
 |B|^2=\frac14|\Ric_0(g)|^2.
 \label{rec:Ricci-beta}
\end{equation}
Set $M=C^2-BB^T$.  Wherever the connection is definite,
\begin{equation}
 \mu_{A_g^+}=\frac23\operatorname{tr}(M)dV_g,\qquad
 Q=\frac{3M}{\operatorname{tr}M}.
 \label{rec:source-Q}
\end{equation}
If $C$ is invertible, the curvature plane is the graph
\begin{equation}
 T=C^{-1}B
 \label{rec:curvature-graph}
\end{equation}
over $\Lambda_g^+$.  In particular, all of these formulas hold in a
sufficiently small $C^2$ neighborhood of $g_0$.
\end{lemma}

\begin{proof}
The curvature operator has blocks
\[
 \mathcal R_g=
 \begin{pmatrix}W^++sI/12&B\\B^T&W^-+sI/12\end{pmatrix}.
\]
To check the normalization of $B$, diagonalize
$g^{-1}\Ric_0=\operatorname{diag}(r_1,r_2,r_3,r_4)$, with
$\sum r_a=0$.  The trace-free Ricci part of the curvature operator has
sectional entries $(r_a+r_b)/2$.  Its off-diagonal chiral entries in
\eqref{rec:form-bases} are therefore
$(r_1+r_2)/2,(r_1+r_3)/2,(r_1+r_4)/2$.  These agree with
\eqref{rec:beta-diagonal}.  Equivariance proves
\eqref{rec:Ricci-beta} in general.

Taking wedge products of \eqref{rec:chiral-curvature} gives
$F_i\wedge F_j=2(C^2-BB^T)_{ij}dV_g$ and hence
\eqref{rec:source-Q}.  Finally, multiplying the column of curvature
forms on the left by $-C^{-1}$ does not change its span and yields
$\Sigma_g^++C^{-1}B\Sigma_g^-$, proving
\eqref{rec:curvature-graph}.  At round, $C=KI$ and $M=K^2I$;
invertibility and positive definiteness persist for small $C^2$
perturbations.
\end{proof}

\subsection{Linear reconstruction at the round metric}

Let $\mathscr R(g)=\widehat g_{A_g^+}$, with the fixed normalization
\eqref{rec:metric-definition}.  Primes below denote derivatives at
$g_0$ along $g_\varepsilon=g_0+\varepsilon h+O(\varepsilon^2)$.
For $\Ric_0'$ we differentiate the trace-free Ricci tensor as a
covariant two-tensor; because $\Ric_0(g_0)=0$, raising one index at
$g_0$ gives the same first-order endomorphism used in the graph map.

\begin{proposition}[The linear reconstruction operator]
\label{rec:linear-reconstruction}
The first variation $k=\mathscr R'_{g_0}h$ is
\begin{equation}
 k=h+\frac1K\Ric_0'+\frac{s'}{12K}g_0
   =\frac1K\Ric'-2h-\frac{s'}{6K}g_0.
 \label{rec:linear-reconstruction-formula}
\end{equation}
On transverse-traceless tensors, with
$L=\nabla^*\nabla=-\nabla^a\nabla_a$,
\begin{equation}
 \Ric_0'=\frac12(L+2K)h,\qquad s'=0,\qquad
 k=\frac{L+4K}{2K}h.
 \label{rec:linear-TT}
\end{equation}
\end{proposition}

\begin{proof}
Write pointwise
\begin{equation}
 \mathscr R(g_\varepsilon)
    =g_\varepsilon(c_\varepsilon e^{H_\varepsilon}\cdot,\cdot),
 \qquad \operatorname{tr}H_\varepsilon=0.
 \label{rec:logmetric-representation}
\end{equation}
The volume condition and \eqref{rec:source-Q} imply
\begin{equation}
 c_\varepsilon^2=\frac{\operatorname{tr}M}{3K^2}.
 \label{rec:scale-c}
\end{equation}
At round, $C=KI$, $B=0$, and $C'=W^{+\prime}+s'I/12$.
Thus $\operatorname{tr}M'=Ks'/2$ and $c'=s'/(12K)$.
The graph formula gives $T'=B'/K$; applying
Lemma~\ref{rec:graph-lemma} and \eqref{rec:Ricci-beta} yields
$H'=g_0^{-1}\Ric_0'/K$.  Differentiating
\eqref{rec:logmetric-representation} proves the first equality in
\eqref{rec:linear-reconstruction-formula}.  The identity
$\Ric_0'=\Ric'-(s'/4)g_0-3Kh$ proves the second.

For the TT formula, the variation of the Christoffel symbols gives
\[
 2\Ric'_{ij}=
 \nabla^a\nabla_i h_{ja}+\nabla^a\nabla_j h_{ia}
             -\nabla^a\nabla_a h_{ij}
             -\nabla_i\nabla_j\operatorname{tr}h.
\]
When $\nabla^a h_{aj}=0$ and $\operatorname{tr}h=0$, commuting
derivatives on the round metric gives
$\nabla^a\nabla_i h_{ja}=4Kh_{ij}$ and the same term with $i,j$
interchanged.  Therefore $\Ric'=\tfrac12(L+8K)h$.
The scalar variation
$s'=\nabla^i\nabla^j h_{ij}-\nabla^a\nabla_a\operatorname{tr}h
      -\langle\Ric,h\rangle$
vanishes.  Subtraction of $3Kh$ proves the formula for $\Ric_0'$,
and substitution proves \eqref{rec:linear-TT}.
\end{proof}

Three checks clarify what reconstruction retains.  Constant rescaling
$h=2cg_0$ gives $k=0$, since the Levi--Civita connection does not change
under a spatially constant scale change.  Naturality gives $k=h$ for
$h=\mathcal L_Xg_0$.  For a conformal perturbation $h=2fg_0$, set
$\Delta f=\nabla^a\nabla_af$.  The conformal Ricci formula gives
\[
 \Ric'=-2\nabla^2f-(\Delta f)g_0,\qquad
 s'=-6\Delta f-24Kf,
\]
so
\begin{equation}
 \mathscr R'_{g_0}(2fg_0)
        =-\frac2K\nabla^2f
        =-\frac1K\mathcal L_{\nabla f}g_0.
 \label{rec:conformal-check}
\end{equation}
Thus conformal source variations disappear modulo diffeomorphisms
at first order.  This does not assert a nonlinear conformal symmetry
of reconstruction or of a metric pullback functional.

On TT tensors, forward reconstruction is an elliptic operator of
order two and amplifies high-frequency modes.  Its inverse on this
sector is $2K(L+4K)^{-1}$.  For comparison, the vector field of
volume-normalized Ricci flow is
$\mathcal F_{\rm RF}(g)=-2\Ric(g)+(\overline s/2)g$ in dimension four.
At round its TT derivative is $-(L+2K)$, since
$\overline s'=0$.  Consequently
\begin{equation}
 \mathscr R'_{g_0}=I-\frac1{2K}\mathcal F_{\rm RF}'
                  \qquad\hbox{on TT tensors}.
 \label{rec:resolvent-comparison}
\end{equation}
This operator identity alone does not identify nonlinear flows.

For later use, denote by
$\mathcal Hh=(A_g^+)'_{g_0}h$ the linearized chiral Levi--Civita map
in the symmetric coframe identification.  Under the same
identifications the normalized triple in
\eqref{rec:adapted-frame} has, for TT data, the first-order variation
of the coframe triple associated with $k=\mathscr R'_{g_0}h$; the
$W^{+\prime}$ contribution is cancelled by $Q^{-1/2}$.
Indeed $Q'=2W^{+\prime}/K$, and the anti-self-dual coefficient of
the triple variation is
$\beta(h)+B'/K=\beta(k)$, with no scalar or skew self-dual term.
Linearizing \eqref{rec:compatibility-equivalence} therefore gives
\begin{equation}
 \mathcal T'=[\mathcal H(h-k),\Sigma_0^+]
     =-K^{-1}[\mathcal H(\Ric_0'),\Sigma_0^+]
                     \qquad\hbox{for TT }h.
 \label{rec:linear-compatibility}
\end{equation}
In particular, belonging to the linearized image of
$g\mapsto A_g^+$ does not imply vanishing compatibility defect with
the reconstructed metric.

\subsection{Second-order reconstruction}

All expansions in this subsection are pointwise in moving
$g_\varepsilon$-orthonormal frames and use coefficients without
factorials.  In particular, the symbols for $\Ric_0$ below denote
its matrix with one index raised; this convention keeps the moving
frame corrections in its coefficients.  Write
\begin{align}
 g_\varepsilon^{-1}\Ric_0(g_\varepsilon)
     &=\varepsilon\rho_1+\varepsilon^2\rho_2+O(\varepsilon^3),
 &B&=\varepsilon B_1+\varepsilon^2B_2+O(\varepsilon^3),
 \label{rec:rho-expansion}\\
 W^+&=\varepsilon W_1+\varepsilon^2W_2+O(\varepsilon^3),
 &s&=12K+\varepsilon s_1+\varepsilon^2s_2+O(\varepsilon^3),
 \label{rec:W-s-expansion}
\end{align}
where $B_j=\beta(\rho_j)$.  The matrices $\rho_j$ and $W_j$ are
symmetric and trace-free, whereas $B_j$ need not be symmetric.

\begin{proposition}[Second-order metric and scale]
\label{rec:second-order}
In \eqref{rec:logmetric-representation}, write
$H_\varepsilon=\varepsilon H_1+\varepsilon^2H_2+O(\varepsilon^3)$.
Then
\begin{align}
 H_1&=\frac{\rho_1}{K},
 \label{rec:H-first}\\
 H_2&=\frac{\rho_2}{K}
       -\frac{s_1\rho_1}{12K^2}
       -\frac{\beta^{-1}(W_1B_1)}{K^2},
 \label{rec:H-second}\\
 c_\varepsilon
   &=1+\frac{\varepsilon s_1}{12K}
       +\varepsilon^2\left(
           \frac{s_2}{12K}
           +\frac{|W_1|^2-|B_1|^2}{6K^2}\right)
       +O(\varepsilon^3).
 \label{rec:c-second}
\end{align}
Equivalently, the relative metric endomorphism is
\begin{align}
 c_\varepsilon e^{H_\varepsilon}
  ={}&I+\varepsilon\left(\frac{\rho_1}{K}
                                +\frac{s_1}{12K}I\right)\notag\\
    &+\varepsilon^2\left[
       \frac{\rho_2}{K}+\frac{s_2}{12K}I
       -\frac{\beta^{-1}(W_1B_1)}{K^2}
       +\frac{\rho_1^2}{2K^2}
       +\frac{|W_1|^2-|B_1|^2}{6K^2}I
                       \right]+O(\varepsilon^3).
 \label{rec:relative-metric-second}
\end{align}
\end{proposition}

\begin{proof}
Put $a_j=s_j/12$.  Expanding the matrix inverse at $C=KI$ gives
\[
 C^{-1}B
  =\frac{\varepsilon}{K}B_1
    +\varepsilon^2\left[
           \frac{B_2}{K}
           -\frac{(W_1+a_1I)B_1}{K^2}\right]
    +O(\varepsilon^3).
\]
By Lemma~\ref{rec:graph-lemma}, the graph of the logarithmic metric
has expansion $\beta(H_\varepsilon)+O(\varepsilon^3)$, with no
quadratic correction.  Applying $\beta^{-1}$ to the first two
coefficients proves \eqref{rec:H-first}--\eqref{rec:H-second}.
The product $W_1B_1$ is an arbitrary graph matrix; its lack of
symmetry causes no problem because $\beta$ maps onto all
three-by-three graph matrices.

For the scale, $\operatorname{tr}W^+=0$ gives the exact identity
\[
 \operatorname{tr}M
     =3(s/12)^2+|W^+|^2-|B|^2.
\]
Inserting the expansions into \eqref{rec:scale-c} and taking the
positive square root proves \eqref{rec:c-second}.  In particular,
the $s_1^2$ term cancels when the square root is expanded.
Finally,
$e^{H_\varepsilon}
 =I+\varepsilon H_1+\varepsilon^2(H_2+H_1^2/2)+O(\varepsilon^3)$.
Multiplication by \eqref{rec:c-second} cancels
$-s_1\rho_1/(12K^2)$ against the product of the first-order scalar
and trace-free terms, giving \eqref{rec:relative-metric-second}.
\end{proof}

The remainder statements are pointwise Taylor statements for a
smooth family; on a compact manifold they are uniform in any fixed
finite norm when the corresponding family and derivatives are
controlled.  They are not estimates along an arbitrary evolution.
In particular, \eqref{rec:relative-metric-second} describes the
nonlinear reconstruction map, not the complete quadratic part of a
connection-induced metric flow.  Determining that flow also requires
the evolution of the independent connection and the compatibility
defect \eqref{rec:torsion-defect}.


\section{The definite-connection functional and its gradient}
\label{sec:definite-flow}

Retain the bundle, definite curvature, and reconstructed metric of Section~\ref{rec:section}. The orientation component is the one containing the round chiral connection. In particular,
\[
F_{0i}=-K\Sigma_i^+,\qquad Q_0=I,\qquad
\mu_0=2K^2dV_{g_0},\qquad \widehat g_{A_0}=g_0.
\]
The constant \(K\) remains fixed.

Consider the energy and the metric on its space of connections
\begin{equation}
 \mathcal P(A)=\int_M\operatorname{tr}\bigl((Q-I)^2\bigr)\mu_A,
 \qquad
 \langle a,b\rangle_A=
 \int_M\sum_i\langle a_i,b_i\rangle_{\widehat g_A}
          dV_{\widehat g_A}.
 \label{flow:energy-and-metric}
\end{equation}
The one-form metric in this formula, including its internal normalization,
is part of the definition of the gradient.  Fine's canonical reconstructed
metric has volume $\mu_A$, and is consequently
$\sqrt2K\,\widehat g_A$.  Our flow differs from the corresponding flow
in~\cite{Fine2011} by this fixed time normalization.

\begin{proposition}
\label{flow:gradient}
With adjoints taken using \eqref{flow:energy-and-metric}, the negative
gradient flow is
\begin{align}
 A_t&=-d_A^*\left[
      \left(4Q-\frac23\operatorname{tr}(Q^2)I\right)F_A\right]
       \label{flow:gradient-equation}\\
    &=-4d_A^*\left[
      \left(q-\frac16\lvert q\rvert^2I\right)F_A\right],
       \qquad q=Q-I.\nonumber
\end{align}
For every smooth solution,
\begin{equation}
 \frac{d}{dt}\mathcal P(A(t))=-\lVert A_t\rVert_A^2.
 \label{flow:energy-dissipation}
\end{equation}
\end{proposition}

\begin{proof}
For a connection variation $a$, write $\dot F=d_Aa$.  Differentiating
\eqref{rec:mu-Q} gives
\begin{equation}
 \dot\mu=\frac23\sum_i(d_Aa)_i\wedge F_i,
 \qquad
 \dot Q_{ij}\mu=(d_Aa)_i\wedge F_j+F_i\wedge(d_Aa)_j
                         -Q_{ij}\dot\mu.
 \label{flow:elementary-variation}
\end{equation}
The Bianchi identity implies
$\dot\mu=\frac23d\sum_i(a_i\wedge F_i)$, so $\int_M\mu_A$ is
constant on the space of connections on this fixed bundle.
Since $\operatorname{tr}Q=3$, we can therefore vary
$\int\operatorname{tr}(Q^2)\mu$ instead of $\mathcal P$.
Using \eqref{flow:elementary-variation}, symmetry of $Q$, and symmetry
of the wedge pairing on two-forms, we obtain
\begin{align}
 \delta\mathcal P(a)
 &=4\int_M\sum_{i,j}Q_{ij}(d_Aa)_i\wedge F_j
       -\frac23\int_M\operatorname{tr}(Q^2)
                     \sum_i(d_Aa)_i\wedge F_i\nonumber\\
 &=\int_M\sum_i(d_Aa)_i\wedge
        \left[\left(4Q-\frac23\operatorname{tr}(Q^2)I\right)F_A\right]_i.
 \label{flow:first-variation}
\end{align}
Each form in the second factor is selfdual for $\widehat g_A$.
Thus the last integral is the ordinary $L^2$ pairing of two-forms,
and integration by parts gives \eqref{flow:gradient-equation}.
There is no additional factor $2K^2$ in this step: the wedge pairing
already incorporates the volume form used in the two-form pairing.

Since $*_AF_A=F_A$ and $d_AF_A=0$, one has $d_A^*F_A=0$.
Writing $Q=I+q$ and $\operatorname{tr}q=0$ in the first line of
\eqref{flow:gradient-equation} now yields the second line.
Finally, pairing the gradient with its negative gives
\eqref{flow:energy-dissipation}.  Dependence of the metric on $A$
does not add a term to this directional derivative of the energy.
\end{proof}

This is Fine's definite-connection gradient equation~\cite{Fine2011},
with the normalization stated above.
The curvature is selfdual for its reconstructed metric at every point of
the configuration space; the equation is consequently not Yang--Mills
heat flow on an externally fixed metric.  Its internal gauge and
diffeomorphism symmetries must be treated when proving parabolicity.
Fine proves short-time existence and uniqueness, and local stability
near perfect connections, by appropriate gauge fixing~\cite{Fine2011}.
These results do not supply a general long-time theorem for arbitrary
initial definite connections.

The exact matrix, volume, and metric velocities for this flow are obtained by inserting its velocity into Proposition~\ref{rec:velocity-proposition}. In particular, equations~\eqref{rec:velocity-Q} and~\eqref{rec:invariant-velocities} retain the trace-normalization term in \(Q_t\); total reconstructed volume is constant while its density may vary.

\section{The closed defect equation at the round connection}
\label{sec:round-defect}

Write
\begin{equation}
 A=A_0+\varepsilon a+O(\varepsilon^2),\qquad
 Q=I+\varepsilon q+O(\varepsilon^2),\qquad
 q\in\Gamma(\operatorname{Sym}^2_0E).
\end{equation}
In this section $q$ denotes the linearized defect, rather than the full
nonlinear matrix $Q-I$.  Because the nonlinear right-hand side of
\eqref{flow:gradient-equation} vanishes at $A_0$, its linearization is
\begin{equation}
 a_t=-4d_0^*(qF_0),\qquad d_0=d_{A_0}.
 \label{flow:linear-connection}
\end{equation}
There is no variation of the adjoint in this expression: in the second
line of \eqref{flow:gradient-equation}, that variation multiplies a
factor which is zero at the background.

We use the selfdual symbols
\begin{equation}
 \eta_{i,jk}=\epsilon_{ijk},\qquad
 \eta_{i,j4}=\delta_{ij},\qquad
 \eta_{i,4j}=-\delta_{ij},\qquad
 \Sigma_i^+=\frac12\eta_{i,ab}e^a\wedge e^b.
 \label{flow:eta-convention}
\end{equation}
They are parallel for the combined tangent and chiral connections.
All covariant derivatives in this section use those round connections.
If $(d_0a)_i=U(a)_{ij}\Sigma_j^++V(a)_{ij}\Sigma_j^-$, then
\begin{equation}
 U(a)_{ik}=\frac14\eta_{k,ab}(d_0a)_{i,ab}
          =\frac12\eta_{k,ab}\nabla_a a_{ib},
 \qquad
 q=-\frac2K\operatorname{sym}_0U(a).
 \label{flow:linear-Q-map}
\end{equation}
Here $\operatorname{sym}_0$ means symmetric trace-free part in the
internal indices.

\begin{proposition}
\label{flow:round-heat}
The linearized defect satisfies
\begin{equation}
 \boxed{q_t=-4(L_E+6K)q,\qquad
 L_E=\nabla_{A_0}^*\nabla_{A_0}
       \quad\hbox{on }\operatorname{Sym}^2_0E.}
 \label{flow:round-defect-heat}
\end{equation}
The linearized local volume velocity vanishes.
\end{proposition}

\begin{proof}
Since $(F_{0j})_{ab}=-K\eta_{j,ab}$, equation
\eqref{flow:linear-connection} is
\begin{equation}
 (a_t)_{ib}=-4K\eta_{j,ab}\nabla_aq_{ij}.
 \label{flow:linear-velocity-components}
\end{equation}
Applying \eqref{flow:linear-Q-map} gives
\begin{equation}
 q_{t,ik}=4\operatorname{sym}_0
  \left[\eta_{k,ab}\eta_{j,cb}\nabla_a\nabla_cq_{ij}\right].
 \label{flow:q-second-derivatives}
\end{equation}
The algebraic identity
\begin{equation}
 \eta_{k,ab}\eta_{j,cb}
   =\delta_{kj}\delta_{ac}+\epsilon_{kjl}\eta_{l,ac}
 \label{flow:eta-product}
\end{equation}
separates this into a rough Laplacian and a curvature term.
The curvature action on the two internal indices is
\begin{equation}
 [\nabla_a,\nabla_c]q_{ij}
   =-K\eta_{t,ac}
      (\epsilon_{itm}q_{mj}+\epsilon_{jtm}q_{im}).
 \label{flow:q-curvature-action}
\end{equation}
Using $\eta_{l,ac}\eta_{t,ac}=4\delta_{lt}$, symmetry of $q$, and
$\operatorname{tr}q=0$, the antisymmetric part of the second derivatives
in \eqref{flow:q-second-derivatives} is
\begin{align}
 \frac12\epsilon_{kjl}\eta_{l,ac}
       [\nabla_a,\nabla_c]q_{ij}
 &=-2K\epsilon_{kjl}
       (\epsilon_{ilm}q_{mj}+\epsilon_{jlm}q_{im})\nonumber\\
 &=-2K\bigl(q_{ki}-\delta_{ik}\operatorname{tr}q+2q_{ik}\bigr)
   =-6Kq_{ik}.
 \label{flow:q-curvature-contraction}
\end{align}
Consequently $q_t=4(\nabla^a\nabla_a-6K)q$, proving the asserted
formula with the nonnegative rough-Laplacian convention.

For the volume assertion, apply \eqref{flow:linear-Q-map} to the
velocity \eqref{flow:linear-velocity-components} and take its trace:
\begin{equation}
 \operatorname{tr}U(a_t)
 =-2K\left[
     \nabla^a\nabla_a\operatorname{tr}q
     +\epsilon_{ijl}\eta_{l,ac}\nabla_a\nabla_cq_{ij}\right]=0.
\end{equation}
The first term vanishes by trace-freeness, and the second by symmetry
in $i,j$.  Equation \eqref{rec:mu-velocity} therefore has no term of
order $\varepsilon$ in its volume velocity.
\end{proof}

The pointwise and integrated consequences are
\begin{align}
 (\partial_t-4\Delta_{g_0})\lvert q\rvert^2
    &=-8\lvert\nabla q\rvert^2-48K\lvert q\rvert^2,
      \label{flow:linear-pointwise-decay}\\
 \frac{d}{dt}\lVert q\rVert_{L^2(g_0)}^2
    &=-8\lVert\nabla q\rVert_{L^2(g_0)}^2
      -48K\lVert q\rVert_{L^2(g_0)}^2.\nonumber
\end{align}
Here $\Delta=\nabla^a\nabla_a$ on functions.  On the compact sphere,
the scalar maximum principle gives
\begin{equation}
 \sup_M\lvert q(t)\rvert
       \leq e^{-24Kt}\sup_M\lvert q(0)\rvert.
 \label{flow:linear-sup-bound}
\end{equation}
These are estimates for the linear equation.  They do not constitute a
nonlinear maximum principle for $Q$, nor do they prove preservation of
a uniform lower bound for its eigenvalues along the nonlinear flow.

\section{The Bianchi identity and the Ricci-flow linearization}
\label{sec:chiral-bianchi}

Let $g$ now be an arbitrary oriented Riemannian four-metric, with its
chiral Levi--Civita connection $A_g^+$.  Local orthonormal frames for
nearby metrics are identified by the symmetric coframe variation
$\dot e^a=\frac12h^a{}_be^b$.  The linearized chiral connection map
$\mathcal H_g h=(A_g^+)'h$ then has the local formula
\begin{equation}
 (\mathcal H_g h)_{ib}
       =\frac12\eta_{i,cd}\nabla_ch_{bd}.
 \label{flow:H-definition}
\end{equation}
Indeed, linearizing the torsion equations gives
$\dot\omega_{ab}=\frac12(\nabla_bh_{ac}-\nabla_ah_{bc})e^c$;
contracting with $-\eta_{i,ab}/2$ gives
\eqref{flow:H-definition}.  This choice of identification fixes the
otherwise available internal gauge term in the variation.

Put $\Ric_0=\operatorname{Ric}-\frac{s}{4}g$.  The Weyl tensor is
viewed either as a four-index tensor or as an endomorphism of two-forms,
with $W^+$ denoting its selfdual block.  Define the $E$-valued two-form
$\Psi=W^+\Sigma^+$.  With the ordinary two-form normalization,
\begin{equation}
 \Psi_{i,ab}=W^+_{ij}\eta_{j,ab}
           =\frac12\eta_{i,cd}W_{abcd}.
 \label{flow:W-projection}
\end{equation}

\begin{lemma}
\label{flow:bianchi-lemma}
For the chiral Levi--Civita connection and adjoint defined by $g$,
\begin{equation}
 \boxed{
 d_{A_g^+}^*(W^+\Sigma^+)
    =-\mathcal H_g(\Ric_0)
      -\frac1{12}\mathcal H_g(sg).
 }
 \label{flow:full-bianchi-identity}
\end{equation}
In this identity the metric is $g$, not the generally different
curvature-reconstructed metric $\widehat g_{A_g^+}$.
\end{lemma}

\begin{proof}
We use $R_{abcd}=K(g_{ac}g_{bd}-g_{ad}g_{bc})$ on the round sphere.
The curvature decomposition in this convention is
\begin{align}
 R_{abcd}=W_{abcd}
  &+\frac12(g_{ac}(\Ric_0)_{bd}+g_{bd}(\Ric_0)_{ac}
              -g_{ad}(\Ric_0)_{bc}-g_{bc}(\Ric_0)_{ad})\nonumber\\
  &+\frac{s}{12}(g_{ac}g_{bd}-g_{ad}g_{bc}).
\end{align}
Contracted Bianchi gives
$\nabla^aR_{abcd}=\nabla_c\operatorname{Ric}_{bd}
                    -\nabla_d\operatorname{Ric}_{bc}$ and
$\nabla^a(\Ric_0)_{ab}=\frac14\nabla_bs$.
Taking the divergence of the decomposition therefore gives
\begin{equation}
 \nabla^aW_{abcd}
 =\frac12(\nabla_c(\Ric_0)_{bd}-\nabla_d(\Ric_0)_{bc})
    +\frac1{24}(g_{bd}\nabla_cs-g_{bc}\nabla_ds).
 \label{flow:W-divergence}
\end{equation}
The symbols $\eta$ are parallel for the combined connections, so
\eqref{flow:W-projection} and \eqref{flow:W-divergence} imply
\begin{align}
 (d_{A_g^+}^*\Psi)_{ib}
 &=-\frac12\eta_{i,cd}\nabla^aW_{abcd}\nonumber\\
 &=-\frac12\eta_{i,cd}\nabla_c(\Ric_0)_{bd}
       -\frac1{24}\eta_{i,cb}\nabla_cs.
\end{align}
The first term is $-\mathcal H_g(\Ric_0)$ by
\eqref{flow:H-definition}; the second is
$-\mathcal H_g(sg)/12$.  This proves the identity, including its
scalar-curvature term, without an Einstein assumption.
\end{proof}

\subsection{The Ricci-flow comparison on transverse-traceless variations}
\label{sec:ricci-comparison}

Return to $g_0$, and write $\mathcal H=\mathcal H_{g_0}$ and
$\mathcal D h=(W^+)'_{g_0}h$.  Set $L=\nabla^*\nabla$ on symmetric tensors.
We restrict to transverse-traceless (TT) tensors,
\begin{equation}
 \nabla^ah_{ab}=0,\qquad \operatorname{tr}_{g_0}h=0.
 \label{flow:TT}
\end{equation}
For these tensors,
\begin{equation}
 s'=0,\qquad
 \operatorname{Ric}'=\frac12(L+8K)h,\qquad
 \Ric_0'=\frac12(L+2K)h.
 \label{flow:linear-Ricci}
\end{equation}
These variation formulas were derived in Proposition~\ref{rec:linear-reconstruction}.

The metric-induced curvature has the block expression
$F=-C\Sigma_g^+-B\Sigma_g^-$, with
$C=W^++sI/12$ and $B$ the trace-free Ricci block.  Consequently,
\begin{equation}
 Q=\frac{3(C^2-BB^T)}{\operatorname{tr}(C^2-BB^T)},\qquad
 \left.\delta Q\right|_{a=\mathcal Hh}=\frac2K\mathcal D h.
 \label{flow:Q-W-intertwining}
\end{equation}
The second identity follows by differentiating at $C=KI$, $B=0$:
the scalar variation cancels under trace normalization, and $BB^T$
has no linear term.

\begin{proposition}
\label{flow:ricci-intertwining}
Let $a=\mathcal Hh$ with $h$ TT.  The linearized independent connection
flow satisfies
\begin{equation}
 \boxed{
 a_t=\mathcal H[-4(L+2K)h].
 }
 \label{flow:ricci-intertwining-equation}
\end{equation}
Thus it preserves the image of TT tensors under $\mathcal H$.
The identity is exact in the stated symmetric coframe identification;
no additional internal gauge term is required there.
\end{proposition}

\begin{proof}
At round, $W^+=0$, $\Ric_0=0$, and $s$ is constant.
Linearizing \eqref{flow:full-bianchi-identity} therefore gives
\begin{equation}
 d_0^*(\mathcal D h\,\Sigma^+)
   =-\mathcal H(\Ric_0')
      -\frac1{12}\mathcal H(s'g_0).
 \label{flow:linear-bianchi}
\end{equation}
In particular, variations of the adjoint and of the triple on the left
multiply the zero Weyl tensor.  On the right,
$\mathcal H_g(sg)$ depends only on $ds$, so the constant background
scalar curvature creates no additional term.
For TT $h$ the scalar term vanishes.
Substituting \eqref{flow:Q-W-intertwining} and $F_0=-K\Sigma^+$ into
\eqref{flow:linear-connection}, and then using
\eqref{flow:linear-Ricci}, gives
\begin{equation}
 a_t=8d_0^*(\mathcal D h\,\Sigma^+)
     =-8\mathcal H(\Ric_0')
     =\mathcal H[-4(L+2K)h].
\end{equation}
On constant curvature, $L$ preserves TT tensors, proving the image
preservation assertion.  This argument uses the full Bianchi identity,
not a projection of the connection flow onto the Levi--Civita image.
\end{proof}

Volume-normalized Ricci flow in dimension four is
$g_t=-2\operatorname{Ric}+\overline s\,g/2$.
Its TT linearization at $g_0$ is $h_t=-(L+2K)h$ by
\eqref{flow:linear-Ricci}.  The metric parameter $h$ in
Proposition~\ref{flow:ricci-intertwining} therefore obeys four times
that equation.  This factor is removed by the time change $\tau=4t$.

The same comparison holds for the reconstructed metric perturbation.
Indeed, differentiating the curvature plane and the volume normalization
at a metric-induced connection gives
\begin{equation}
 \widehat h=h+\frac1K\Ric_0'+\frac{s'}{12K}g_0,
 \qquad
 \widehat h=\frac{L+4K}{2K}h\quad\hbox{on TT tensors}.
 \label{flow:reconstruction-linear}
\end{equation}
This is the reconstruction formula of Proposition~\ref{rec:linear-reconstruction}.
The TT reconstruction operator is a polynomial in $L$, so it commutes
with the generator in \eqref{flow:ricci-intertwining-equation} and yields
\begin{equation}
 \boxed{\widehat h_t=-4(L+2K)\widehat h.}
 \label{flow:reconstructed-Ricci}
\end{equation}
This assertion concerns the reconstructed perturbations arising from
the metric-induced TT data just considered.

For a general metric variation, the same Bianchi calculation gives
\begin{equation}
 a_t=\mathcal H\left[-8\Ric_0'-\frac23s'g_0\right].
 \label{flow:general-linear-metric-image}
\end{equation}
In particular, for $h=2fg_0$ one has $\mathcal D h=0$, hence $q=a_t=0$.
Thus the comparison with normalized Ricci flow is a TT statement,
and does not identify their general conformal variations.
Changing the identification of spin frames changes the component
connection variation by an infinitesimal internal gauge transformation;
the statements above specify the identification in which the operator
identity holds.  Diffeomorphism and internal gauge directions lie in the
kernel of the linearized defect at $Q=I$.

\subsection{An independent norm identity for the Levi--Civita variation}

\begin{proposition}
\label{flow:H-norm}
Using $g_0$ for both the tensor and connection $L^2$ pairings and the
internal metric $\delta_{ij}$, one has
\begin{equation}
 \mathcal H^*\mathcal Hh=\frac14(L+4K)h
       \quad\hbox{for every TT tensor }h.
 \label{flow:H-norm-equation}
\end{equation}
In particular, $\mathcal H$ is injective on TT tensors on the round
sphere.
\end{proposition}

\begin{proof}
Let $h$ be TT and let $k$ be any smooth symmetric tensor.  The identity
\begin{equation}
 \sum_i\eta_{i,ab}\eta_{i,de}
      =\delta_{ad}\delta_{be}-\delta_{ae}\delta_{bd}
          +\epsilon_{abde}
\end{equation}
and \eqref{flow:H-definition} give
\begin{align}
 4\langle\mathcal Hh,\mathcal Hk\rangle_{L^2}
  =\int_M\bigl(&\nabla_ah_{cb}\nabla_ak_{cb}
           -\nabla_ah_{cb}\nabla_bk_{ca}\nonumber\\
           &+\epsilon_{abde}\nabla_ah_{cb}\nabla_dk_{ce}\bigr)dV_{g_0}.
 \label{flow:H-bilinear}
\end{align}
The mixed derivative integral is
\begin{align}
 \int_M\nabla_ah_{cb}\nabla_bk_{ca}\,dV
   &=-\int_Mk_{ca}\nabla^b\nabla_ah_{cb}\,dV\nonumber\\
   &=-\int_Mk_{ca}[\nabla^b,\nabla_a]h_{cb}\,dV
    =-4K\int_M\langle h,k\rangle\,dV.
 \label{flow:H-mixed-term}
\end{align}
Here the TT conditions and the constant-curvature action give
$[\nabla^b,\nabla_a]h_{cb}=4Kh_{ca}$.
For the epsilon term, integrate in $d$ and use its antisymmetry in
$a,d$:
\begin{equation}
 \int_M\epsilon_{abde}\nabla_ah_{cb}\nabla_dk_{ce}\,dV
   =-\frac12\int_M\epsilon_{abde}k_{ce}
                         [\nabla_d,\nabla_a]h_{cb}\,dV=0.
 \label{flow:H-epsilon-term}
\end{equation}
To verify the last equality explicitly, the commutator is
\begin{equation}
 [\nabla_d,\nabla_a]h_{cb}
   =-K\bigl(\delta_{ca}h_{db}-\delta_{cd}h_{ab}
               +\delta_{ba}h_{cd}-\delta_{bd}h_{ca}\bigr).
\end{equation}
In the first two terms, the epsilon tensor contracts a symmetric pair
of indices of $h$; in the last two it has a repeated index.
Each contribution is zero.  Finally, the first term in
\eqref{flow:H-bilinear} integrates to $\langle Lh,k\rangle$.
Combining these identities proves
$4\langle\mathcal Hh,\mathcal Hk\rangle
 =\langle(L+4K)h,k\rangle$ for every symmetric $k$, and hence the
claimed operator identity.  Positivity of $K$ gives injectivity.
\end{proof}

This identity also checks the time normalization without using a
Weyl-energy Hessian formula.  Let $\mathscr B a=\delta Q(a)$ and give
its target the pairing $\int\operatorname{tr}(q_1q_2)\mu_0$.
We identify each Hessian with its self-adjoint operator using the stated
configuration-space $L^2$ metric when it is applied to a single argument.
At the perfect connection,
\begin{equation}
 \operatorname{Hess}_{A_0}\mathcal P(a,b)
      =2\int_M\operatorname{tr}(\mathscr B a\,\mathscr B b)\mu_0,
 \qquad
 \operatorname{Hess}_{A_0}\mathcal P(\mathcal Hh)
      =4\mathcal H(L+2K)h.
 \label{flow:connection-hessian}
\end{equation}
The second formula follows from the independently established negative
gradient linearization in Proposition~\ref{flow:ricci-intertwining}.
The metric pullback Hessian will be derived from these identities in Section~\ref{sec:metric-hessian}. This derivation separates the choice of configuration-space inner product from the energy itself. No linear intertwining identity here asserts preservation of the nonlinear Levi--Civita image or equality with nonlinear Ricci flow.


\section{Metric energies and their first nonlinear discrepancy}
\label{sec:metric-energy}

Fix the round metric $g_0$ on $S^4$, with sectional curvature $K>0$.
All Weyl norms in this section are Hilbert--Schmidt norms as curvature
operators on two-forms. In particular, the squared four-index tensor norm
is four times the squared operator norm. We use the metric $L^2$ pairing
on symmetric two-tensors unless a connection pairing is explicitly named.

For a metric $g$ sufficiently close to $g_0$, write the curvature of its
positive chiral Levi--Civita connection in orthogonal chiral frames as
\begin{equation}\label{metric:curvature-blocks}
 F_i=-C_{ij}\Sigma_j^+-B_{ij}\Sigma_j^- ,\qquad
 C=aI+W,\qquad a=\frac{s_g}{12},\qquad W=W_g^+.
\end{equation}
Here $B$ is the off-diagonal curvature block, encoding trace-free Ricci
curvature, and
$\Sigma_i^\pm\wedge\Sigma_j^\pm=\pm2\delta_{ij}\,dV_g$.
Define
\begin{equation}\label{metric:M-Q}
 M=C^2-BB^T,\qquad t=\operatorname{tr}M,\qquad
 \mu_A=\frac23t\,dV_g,\qquad Q=\frac{3M}{t}.
\end{equation}
The connection is definite on the neighborhood under consideration.
At $g_0$, one has $a=K$, $W=B=0$, $M=K^2I$, and $Q=I$.

We distinguish the connection functional and its pullback to metrics:
\begin{equation}\label{metric:functionals}
 \mathcal P_{\mathrm{conn}}(A)=\mathcal P(A),
 \qquad
 \mathcal P_{\mathrm{met}}(g)=\mathcal P_{\mathrm{conn}}(A_g^+),
 \qquad
 \mathcal E_+(g)=\int_{S^4}\|W_g^+\|^2dV_g.
\end{equation}
The two appearances of $\mathcal P$ have different gradient flows because
their domains and $L^2$ pairings differ. Pulling back a functional is not
the same operation as pushing forward its gradient vector field.

\subsection{An exact algebraic formula}

For an internal matrix $T$, let $T_0=T-\tfrac13\operatorname{tr}(T)I$.

\begin{proposition}\label{prop:metric-exact-energy}
On the definite neighborhood of the round metric,
\begin{equation}\label{metric:exact-energy}
 \mathcal P_{\mathrm{met}}(g)
 =6\int_{S^4}
 \frac{\|[2aW+W^2-BB^T]_0\|^2}
 {\operatorname{tr}[(aI+W)^2-BB^T]}\,dV_g.
\end{equation}
More explicitly, set
\[
 w_2=\operatorname{tr}W^2,\qquad
 b_2=\operatorname{tr}BB^T,\qquad
 N_2=(W^2-BB^T)_0,
\]
and
\[
 C_3=\operatorname{tr}W^3-\operatorname{tr}(WBB^T).
\]
Then $t=3a^2+w_2-b_2$, and
\begin{equation}\label{metric:exact-discrepancy}
 \frac18\mathcal P_{\mathrm{met}}(g)-\mathcal E_+(g)
 =\int_{S^4}
 \frac{3aC_3+\tfrac34\|N_2\|^2-w_2^2+b_2w_2}{t}\,dV_g.
\end{equation}
\end{proposition}

\begin{proof}
The wedge identities give
$F_i\wedge F_j=2M_{ij}\,dV_g$. Since
$Q-I=3M_0/t$ and $\mu_A=2t\,dV_g/3$, the integrand in
$\mathcal P_{\mathrm{met}}$ is $6\|M_0\|^2/t$.
Now $M_0=2aW+N_2$, proving \eqref{metric:exact-energy}.
Expanding its numerator gives
\[
 \|2aW+N_2\|^2=4a^2w_2+4aC_3+\|N_2\|^2.
\]
After division by eight and subtraction of $w_2$, the first term becomes
$(3a^2-t)w_2/t=(-w_2^2+b_2w_2)/t$. This proves
\eqref{metric:exact-discrepancy}.
\end{proof}

For example, if $|a-K|\leq K/2$ and $t\geq cK^2$ for a fixed
$c>0$, the exact identity implies
\begin{equation}\label{metric:discrepancy-bound}
 \left|\frac18\mathcal P_{\mathrm{met}}-\mathcal E_+\right|
 \leq \frac{C_c}{K}\int |W|(|W|^2+|B|^2)\,dV_g
 +\frac{C_c}{K^2}\int (|W|^2+|B|^2)^2\,dV_g.
\end{equation}
This is a local curvature estimate with an explicit denominator
condition. It does not control the discrepancy from the Weyl $L^2$
energy alone.

\subsection{Second-order defects and the cubic energy difference}

Consider a smooth metric path $g_\varepsilon=g_0+\varepsilon h+
O(\varepsilon^2)$. Use moving orthonormal frames and write
\begin{align}
 W&=\varepsilon W_1+\varepsilon^2W_2+O(\varepsilon^3),&
 a&=K+\varepsilon a_1+\varepsilon^2a_2+O(\varepsilon^3),
 \label{metric:expansions}\\
 B&=\varepsilon B_1+\varepsilon^2B_2+O(\varepsilon^3),&
 dV_{g_\varepsilon}&=(1+\varepsilon v_1+O(\varepsilon^2))dV_{g_0}.
 \nonumber
\end{align}
These are coefficient expansions, without factorials; in particular
$v_1=\tfrac12\operatorname{tr}_{g_0}h$.

\begin{proposition}\label{prop:metric-cubic}
The normalized curvature matrix has expansion
\begin{equation}\label{metric:Q-second-order}
 Q-I=\frac{2\varepsilon}{K}W_1
 +\varepsilon^2\left[
 \frac2K W_2-\frac{2a_1}{K^2}W_1
 +\frac{(W_1^2-B_1B_1^T)_0}{K^2}\right]
 +O(\varepsilon^3).
\end{equation}
Moreover,
\begin{equation}\label{metric:cubic-discrepancy}
 \frac18\mathcal P_{\mathrm{met}}(g_\varepsilon)
 -\mathcal E_+(g_\varepsilon)
 =\frac{\varepsilon^3}{K}\int_{S^4}
 \left[\operatorname{tr}W_1^3
 -\operatorname{tr}(W_1B_1B_1^T)\right]dV_{g_0}
 +O(\varepsilon^4).
\end{equation}
In particular the cubic coefficient is independent of the acceleration
of the chosen metric path.
\end{proposition}

\begin{proof}
Expansion of $M$ gives
\begin{align*}
 M_0={}&2K\varepsilon W_1
 +\varepsilon^2\left[2KW_2+2a_1W_1
 +(W_1^2-B_1B_1^T)_0\right]+O(\varepsilon^3),\\
 t={}&3K^2+6Ka_1\varepsilon+O(\varepsilon^2).
\end{align*}
Multiplying $M_0$ by $3/t$ proves \eqref{metric:Q-second-order}.
For the energy calculation,
\[
 \mu_A=2K^2\left[1+\varepsilon\left(\frac{2a_1}{K}+v_1\right)
 +O(\varepsilon^2)\right]dV_{g_0}.
\]
Consequently
\begin{align*}
 \mathcal P_{\mathrm{met}}(g_\varepsilon)
 ={}&8\varepsilon^2\int |W_1|^2dV_{g_0}\\
 &+8\varepsilon^3\int\left[
 2\langle W_1,W_2\rangle+v_1|W_1|^2
 +K^{-1}\langle W_1,(W_1^2-B_1B_1^T)_0\rangle
 \right]dV_{g_0}+O(\varepsilon^4).
\end{align*}
The terms containing $a_1$ have canceled.
The corresponding expansion of $\mathcal E_+$ consists of the first two
terms inside the brackets, with the same quadratic coefficient after
division by eight. This proves \eqref{metric:cubic-discrepancy}.
The trace-free property of $W_1$ removes the subscript zero in its final
pairing. These traces are unchanged by the orthogonal changes of chiral
frames. Neither $W_2$ nor $v_1$ occurs in the difference.
\end{proof}

Let
\begin{equation}\label{metric:D-J-definitions}
 \mathcal D h=(W^+)'_{g_0}h,\qquad
 \mathcal J h=B'_{g_0}h.
\end{equation}
Both are second-order differential operators. Their formal adjoints use
the background metric pairing and Hilbert--Schmidt matrix pairings.

\begin{corollary}\label{cor:metric-gradient-discrepancy}
Under any smooth identification of the varying tensor bundles with the
background bundles, the metric gradients satisfy
\begin{align}
 &\operatorname{grad}\left(\frac18\mathcal P_{\mathrm{met}}
 -\mathcal E_+\right)(g_0+\varepsilon h+O(\varepsilon^2))
 \label{metric:gradient-discrepancy}\\
 &\quad=\frac{\varepsilon^2}{K}
 \left\{\mathcal D^*\left[3(W_1^2)_0-(B_1B_1^T)_0\right]
 -2\mathcal J^*(W_1B_1)\right\}+O(\varepsilon^3),
 \nonumber
\end{align}
where $W_1=\mathcal Dh$ and $B_1=\mathcal Jh$.
\end{corollary}

\begin{proof}
The difference of the functionals has vanishing Taylor coefficients
through degree two. Its homogeneous cubic coefficient is
\[
 \mathcal C(h)=\frac1K\int
 \left[\operatorname{tr}(\mathcal Dh)^3
 -\operatorname{tr}\left((\mathcal Dh)(\mathcal Jh)(\mathcal Jh)^T\right)
 \right]dV_{g_0}.
\]
Differentiating in a direction $v$ gives
\[
 d\mathcal C_h(v)=\frac1K\int
 \left\langle3(W_1^2)_0-(B_1B_1^T)_0,\mathcal Dv\right\rangle
 -2\left\langle W_1B_1,\mathcal Jv\right\rangle\,dV_{g_0}.
\]
Integration by parts proves the formula. Since this gradient begins in
degree two, variation of its $L^2$ pairing contributes only in degree
three and higher. The two metric flows thus agree to linear order and
first differ quadratically. Each term in their difference is a nonlinear
fourth-order differential expression; the product of two curvature
variations does not make the equation sixth order.
\end{proof}

\section{A self-contained factorization of the round metric Hessian}
\label{sec:metric-hessian}

We derive the round Weyl-energy factorization from the connection
identities proved above. In particular, it is not used as an input to
the connection calculation. This also checks the relative normalization
of the two configuration spaces.

Write $\mathcal Hh=(A_g^+)'_{g_0}h$. In the symmetric-frame
identification its components are
\begin{equation}\label{metric:H-definition}
 (\mathcal Hh)_{ib}=\frac12\eta_{i,cd}\nabla_ch_{bd}.
\end{equation}
Let $L=\nabla^*\nabla$ on symmetric tensors. For transverse-traceless
tensors, the integrated connection-variation identity
\eqref{flow:H-norm-equation} and the Bianchi intertwining identity
\eqref{flow:ricci-intertwining-equation} state, in the present normalizations,
\begin{equation}\label{metric:H-identities}
 \mathcal H^*\mathcal H=\frac14(L+4K),\qquad
 \mathcal L_{\mathrm{conn}}\mathcal Hh
 =4\mathcal H(L+2K)h.
\end{equation}
Here $\mathcal L_{\mathrm{conn}}$ is the Hessian of
$\mathcal P_{\mathrm{conn}}$ at the round connection for the connection
pairing based on the reconstructed metric normalized to equal $g_0$.
These identities follow from variation of torsion,
the curvature Bianchi identity, and integration by parts.

\begin{theorem}\label{thm:metric-Hessian-factorization}
For every transverse-traceless symmetric tensor $h$ on round $S^4$,
\begin{equation}\label{metric:D-factorization}
 \mathcal D^*\mathcal D h
 =\frac1{16}(L+2K)(L+4K)h.
\end{equation}
Equivalently, the metric Hessians are
\begin{equation}\label{metric:Hessian-pair}
 \mathcal L_{\mathcal P_{\mathrm{met}}}
 =(L+2K)(L+4K),\qquad
 \mathcal L_{\mathcal E_+}
 =\frac18(L+2K)(L+4K)
 \quad\text{on TT tensors}.
\end{equation}
\end{theorem}

\begin{proof}
At the round connection, $Q=I$, so
$d\mathcal P_{\mathrm{conn}}=0$. The ordinary Hessian chain rule for
$\mathcal P_{\mathrm{met}}=\mathcal P_{\mathrm{conn}}\circ(g\mapsto A_g^+)$
therefore has no term involving the second derivative of $g\mapsto A_g^+$:
\[
 \operatorname{Hess}_{g_0}\mathcal P_{\mathrm{met}}(h,k)
 =\left\langle\mathcal Hh,
 \mathcal L_{\mathrm{conn}}\mathcal Hk\right\rangle_{L^2}.
\]
Use \eqref{metric:H-identities}, noting that $L$ preserves TT tensors on
the round sphere, to obtain
\begin{align*}
 \mathcal L_{\mathcal P_{\mathrm{met}}}h
 &=4\mathcal H^*\mathcal H(L+2K)h\\
 &=(L+4K)(L+2K)h.
\end{align*}
On the other hand, Proposition~\ref{prop:metric-cubic} gives
\[
 \operatorname{Hess}_{g_0}\mathcal P_{\mathrm{met}}(h,k)
 =16\langle\mathcal Dh,\mathcal Dk\rangle_{L^2},
 \qquad
 \operatorname{Hess}_{g_0}\mathcal E_+(h,k)
 =2\langle\mathcal Dh,\mathcal Dk\rangle_{L^2}.
\]
Comparison proves the stated identities on the TT subspace. Since
$\mathcal D$ annihilates infinitesimal conformal and diffeomorphism
variations, $\mathcal D^*\mathcal D$ takes values in TT tensors, so the
bilinear identity also gives the operator identity in
\eqref{metric:D-factorization}.
\end{proof}

The factorization agrees with the independently established
quadratic-curvature Hessian formula of
Gursky and Viaclovsky~\cite{GurskyViaclovsky2015}.
The proof above uses the connection and Bianchi identities already
derived in this manuscript, and does not invoke that formula.

The linear negative metric-gradient equations are thus
\begin{equation}\label{metric:linear-flows}
 h_t^{(\mathcal P_{\mathrm{met}})}=-(L+2K)(L+4K)h,
 \qquad
 h_t^{(\mathcal P_{\mathrm{met}}/8)}
 =h_t^{(\mathcal E_+)}=-\frac18(L+2K)(L+4K)h.
\end{equation}
The positivity of both factors gives coercivity on TT tensors and shows
that the kernel of $\mathcal D$ consists precisely of infinitesimal
diffeomorphism and conformal variations, after the standard orthogonal
decomposition of symmetric tensors. In particular a nonzero tensor in
the image of $\mathcal D^*$ cannot lie in $\ker\mathcal D$.

For the connection functional with the stated normalization, the independent
connection flow instead has a second-order linear equation for
$q=Q-I$. Its time variable has length-squared scaling, whereas the
metric equations \eqref{metric:linear-flows} have length-to-the-fourth
scaling. Their numerical rates are not directly comparable without a
choice of time normalization.

\subsection{A TT spectral gap and explicit linear decay}
\label{subsec:metric-gap}

The following estimate needs no tabulation of spherical harmonics.

\begin{lemma}\label{lem:metric-TT-gap}
For every smooth TT tensor $h$ on the round four-sphere,
\begin{equation}\label{metric:TT-gap}
 \langle Lh,h\rangle_{L^2}=\|\nabla h\|_{L^2}^2
 \geq 8K\|h\|_{L^2}^2.
\end{equation}
\end{lemma}

\begin{proof}
Define the fully symmetrized covariant derivative, without averaging, by
\[
 S_{abc}=\nabla_a h_{bc}+\nabla_b h_{ca}+\nabla_c h_{ab}.
\]
Symmetry of $h$ and renaming the summed indices give
\[
 \|S\|_{L^2}^2=3\|\nabla h\|_{L^2}^2+6I,
 \qquad I=\int (\nabla_a h_{bc})(\nabla^b h^{ca})\,dV_{g_0}.
\]
Integrating the first derivative by parts and using
$\nabla^a h_{ca}=0$ yields
\[
 I=-\int h^{bc}\nabla^a\nabla_b h_{ca}\,dV_{g_0}
   =-\int h^{bc}[\nabla^a,\nabla_b]h_{ca}\,dV_{g_0}.
\]
The round curvature commutator on a symmetric tensor is
\[
 [\nabla^a,\nabla_b]h_{ca}
 =4Kh_{bc}-K(\operatorname{tr}h)g_{bc}.
\]
As $h$ is trace-free, $I=-4K\|h\|_{L^2}^2$. Consequently
\begin{equation}\label{metric:symmetrized-derivative-identity}
 \|S\|_{L^2}^2
 =3\|\nabla h\|_{L^2}^2-24K\|h\|_{L^2}^2.
\end{equation}
Nonnegativity of the left-hand side proves the claim.
\end{proof}

\begin{corollary}\label{cor:metric-linear-decay}
For the linearized negative metric gradient flow of either
$\mathcal E_+$ or $\mathcal P_{\mathrm{met}}/8$, a TT rough-Laplacian
eigenmode with eigenvalue $\lambda$ evolves by
\begin{equation}\label{metric:mode-decay}
 h(t)=\exp\!\left[-\frac{(\lambda+2K)(\lambda+4K)}8t\right]h(0).
\end{equation}
Every TT solution satisfies
\begin{equation}\label{metric:TT-decay}
 \|h(t)\|_{L^2}\leq e^{-15K^2t}\|h(0)\|_{L^2}.
\end{equation}
For the independent connection flow, restricted to $a=\mathcal Hh$
with $h$ TT, the corresponding bounds are
\begin{equation}\label{metric:connection-TT-decay}
 \|h(t)\|_{L^2}\leq e^{-40Kt}\|h(0)\|_{L^2},\qquad
 \|a(t)\|_{L^2}\leq e^{-40Kt}\|a(0)\|_{L^2}.
\end{equation}
These are bounds for the linear equations; no assertion that the lower
bound $\lambda=8K$ is attained is needed.
\end{corollary}

\begin{proof}
Equation~\eqref{metric:linear-flows} gives
\eqref{metric:mode-decay}. Since $L$ is self-adjoint and preserves TT
tensors, its orthonormal TT eigenmode expansion and
Lemma~\ref{lem:metric-TT-gap} give $\lambda\geq8K$ on every mode.
Thus $(\lambda+2K)(\lambda+4K)/8\geq15K^2$, proving
\eqref{metric:TT-decay}. The independent connection equation
\eqref{flow:ricci-intertwining-equation} gives
$h_t=-4(L+2K)h$ in this parametrization. Its modal rate is
$4(\lambda+2K)\geq40K$. Finally,
\eqref{metric:H-identities} gives
\[
 \|\mathcal Hh\|_{L^2}^2
 =\frac14\langle h,(L+4K)h\rangle_{L^2};
\]
the same exponential bound therefore holds with the positive mode
weights $(\lambda+4K)/4$, proving the bound for $a$.
\end{proof}

\subsection{Conformal coordinates for the quadratic Weyl energy}
\label{subsec:metric-stereographic}

There is also a direct conformal-coordinate simplification. In a
stereographic chart write
\[
 g_0=\Omega^2\delta,\qquad
 \Omega(x)=\frac{2}{1+K|x|^2},
\]
and express a compactly supported symmetric variation as
$h=\Omega^2 k$. Here $k$ is a symmetric two-tensor on the Euclidean
chart, not a scalar conformal factor. Let
$\mathcal D_{\mathrm{flat}}k=(W^+)'_\delta k$.

\begin{proposition}\label{prop:metric-conformal-covariance}
In corresponding oriented orthonormal chiral frames,
\begin{equation}\label{metric:conformal-D}
 \mathcal D_{g_0}(\Omega^2k)=\Omega^{-2}\mathcal D_{\mathrm{flat}}k.
\end{equation}
Consequently the quadratic Weyl energy satisfies
\begin{equation}\label{metric:flat-quadratic-energy}
 \int_{S^4}|\mathcal D_{g_0}h|^2dV_{g_0}
 =\int_{\mathbb R^4}|\mathcal D_{\mathrm{flat}}k|^2dx.
\end{equation}
\end{proposition}

\begin{proof}
For a fixed conformal factor $\Omega$, Weyl curvature as a curvature
operator transforms by
$W^+_{\Omega^2g}=\Omega^{-2}W^+_g$ in the corresponding orthonormal
frames. The self-dual splitting is conformally invariant in four
dimensions. Apply this identity to
$g_\varepsilon=\Omega^2(\delta+\varepsilon k)$ and differentiate at
$\varepsilon=0$. Since the background Weyl curvature vanishes, changes
of frames introduce no additional term. This proves
\eqref{metric:conformal-D}. The squared norm contributes
$\Omega^{-4}$ and $dV_{g_0}=\Omega^4dx$, giving
\eqref{metric:flat-quadratic-energy}. Compact support makes extension
by zero across the missing stereographic pole smooth.
\end{proof}

Thus conformal coordinates remove the background conformal factor from
this quadratic curvature energy. They retain the tensorial degrees of
freedom of $k$; in particular the TT conditions on $h$ need not become
Euclidean TT conditions on $k$. The identity alone does not reduce the
metric evolution to a scalar equation.

\section{An analytical distinction between the metric flows}
\label{subsec:metric-conformal-test}

The cubic discrepancy already distinguishes the nonlinear metric
gradient vector fields. The following explicit example requires no
numerical evaluation.

\begin{proposition}\label{prop:metric-nonpreservation}
There are real-analytic conformal metrics arbitrarily close to the round
metric for which the negative metric $L^2$-gradient of
$\mathcal P_{\mathrm{met}}/8$ is not tangent to the conformally flat
locus. The Weyl-energy gradient vanishes at all of these metrics.
\end{proposition}

\begin{proof}
Realize the radius-$\ell=K^{-1/2}$ sphere in $\mathbb R^5$, choose its north pole
$p$, and let $X^1,X^2$ be the first two ambient coordinate functions.
Set
\[
 f=\frac12(X^1)^2X^2,\qquad g_\varepsilon=e^{2\varepsilon f}g_0.
\]
This is a globally defined real-analytic family and is conformally flat.
For sufficiently small $\varepsilon$, its chiral connection remains
definite by openness on the compact sphere. In round normal coordinates
$y$ at $p$,
\[
 f=\frac12(y^1)^2y^2+O(|y|^5).
\]
The first variation of trace-free Ricci curvature is
$-2(\nabla^2f)_0$. In the chiral bases
\[
 \frac1{\sqrt2}(e^{12}\pm e^{34},\ e^{13}\mp e^{24},\ e^{14}\pm e^{23}),
\]
its off-diagonal block has the spatial jet
\begin{equation}\label{metric:conformal-B-jet}
 B_1(y)=y^1S-\frac{y^2}{2}I+O(|y|^3),\qquad
 S=\begin{pmatrix}0&0&0\\0&0&1\\0&-1&0\end{pmatrix}.
\end{equation}
For completeness, the Hessian of the displayed cubic at leading order is
\[
 \nabla^2f=
 \begin{pmatrix}y^2&y^1&0&0\\y^1&0&0&0\\0&0&0&0\\0&0&0&0\end{pmatrix}
 +O(|y|^3).
\]
Using $\operatorname{Ric}_0'=-2(\nabla^2 f)_0$ in the curvature decomposition and projecting to the
opposite chiral blocks gives \eqref{metric:conformal-B-jet}.
As $S^T=-S$,
\begin{equation}\label{metric:conformal-N-jet}
 N:=(B_1B_1^T)_0
 =(y^1)^2T+O(|y|^4),\qquad
 T=\operatorname{diag}(-2/3,1/3,1/3).
\end{equation}
In particular $N(p)=\nabla N(p)=0$, so only the flat principal part of
$\mathcal D^*$ contributes to $\mathcal D^*N(p)$.

We compute that principal part explicitly. For a flat variation $k$,
\[
 R'(k)_{ijkl}=\frac12\left(
 \partial_j\partial_k k_{il}+\partial_i\partial_l k_{jk}
 -\partial_j\partial_l k_{ik}-\partial_i\partial_k k_{jl}\right).
\]
Remove its Ricci and scalar traces and take its positive chiral block.
Let $E^{aa}$ denote the diagonal tensor with its sole nonzero entry in
position $(a,a)$, and set $\phi=(y^1)^2/2$. The four results are
\begin{align*}
 \mathcal D_{\mathrm{flat}}(E^{11}\phi)&=0,\\
 \mathcal D_{\mathrm{flat}}(E^{22}\phi)
 &=\operatorname{diag}(-1/6,1/12,1/12),\\
 \mathcal D_{\mathrm{flat}}(E^{33}\phi)
 &=\operatorname{diag}(1/12,-1/6,1/12),\\
 \mathcal D_{\mathrm{flat}}(E^{44}\phi)
 &=\operatorname{diag}(1/12,1/12,-1/6).
\end{align*}
The corresponding off-diagonal symmetric tensors give matrices whose
diagonal entries vanish. Since $\partial_1^2N(p)=2T$ and all other
second derivatives of $N$ vanish at $p$, the adjoint pairing therefore
gives
\begin{equation}\label{metric:nonzero-transverse-source}
 \mathcal D^*N(p)=\operatorname{diag}(0,1/3,-1/6,-1/6)\ne0.
\end{equation}

For the conformal family, $W_1=0$ and the full Weyl energy vanishes
identically. Corollary~\ref{cor:metric-gradient-discrepancy} consequently
gives
\begin{equation}\label{metric:conformal-velocity}
 -\operatorname{grad}(\mathcal P_{\mathrm{met}}/8)(g_\varepsilon)
 =\frac{\varepsilon^2}{K}\mathcal D^*N+O(\varepsilon^3).
\end{equation}
The global tensor $\mathcal D^*N$ is TT because $\mathcal D$
annihilates conformal and diffeomorphism directions. It is nonzero by
\eqref{metric:nonzero-transverse-source}. The adjoint identity
\[
 0<\|\mathcal D^*N\|_{L^2}^2
 =\langle N,\mathcal D\mathcal D^*N\rangle_{L^2}
\]
therefore shows that $\mathcal D\mathcal D^*N\ne0$, independently of the
Hessian factorization. Applying the linearization of $W^+$ at
$g_\varepsilon$ to \eqref{metric:conformal-velocity} thus gives a
nonzero result for all sufficiently small nonzero $\varepsilon$.
The gradient vector field is not tangent to $W^+=0$, and hence is not
tangent to the conformally flat locus. Meanwhile
$\operatorname{grad}\mathcal E_+=0$ at a metric with $W^+=0$.
\end{proof}

There is no conflict with the quartic onset of the connection energy
along a conformal family. For any such family,
\begin{equation}\label{metric:conformal-quartic-energy}
 \mathcal P_{\mathrm{met}}(g_\varepsilon)
 =\frac{2\varepsilon^4}{K^2}\int\|(B_1B_1^T)_0\|^2dV_{g_0}
 +O(\varepsilon^5).
\end{equation}
The derivative along the conformal family tests only conformal
directions; the quadratic leading gradient in
\eqref{metric:conformal-velocity} is transverse to them. This proposition
concerns the metric pullback functional and its gradient. It makes no
nonpreservation assertion about an independent connection flow with its
different pairing and reconstruction map.


\section{Energy bounds and the boundary of definiteness}\label{sec:boundary}

The defect functional measures the failure of a definite connection to be perfect. Its algebraic normalization also limits what estimates can follow from energy monotonicity.

\begin{proposition}[The defect is bounded on the definite cone]\label{prop:bounded-defect}
If \(Q>0\) and \(\operatorname{tr}Q=3\), then
\[
0\leq |Q-I|^2<6.
\]
For a connection on the fixed round chiral bundle,
\[
0\leq\mathcal P(A)<32\pi^2.
\]
If \(m=\lambda_{\min}(Q)\), the pointwise estimate
\begin{equation}\label{eq:min-eigen-defect}
m\geq1-\sqrt{\frac23}\,|Q-I|
\end{equation}
holds. In particular, a sufficiently small \emph{uniform} defect yields a positive eigenvalue bound; a finite upper bound on the integrated energy does not provide such a conclusion.
\end{proposition}
\begin{proof}
Write the positive eigenvalues as \(\lambda_1,\lambda_2,\lambda_3\), with sum \(3\). Then
\[
|Q-I|^2=\sum_i\lambda_i^2-3<9-3.
\]
Integrating against \(\mu_A\) and using
\(\int\mu_A=16\pi^2/3\) gives the energy bound. Since \(m\leq1\), minimizing the other two squared deviations at fixed sum gives
\[
|Q-I|^2\geq(m-1)^2+
2\left(\frac{3-m}{2}-1\right)^2
=\frac32(1-m)^2.
\]
This is \eqref{eq:min-eigen-defect}. Neither inequality controls the spatial concentration of the defect.
\end{proof}

\begin{proposition}[Finite energy density at a degenerate reconstruction]\label{prop:boundary-example}
There is a smooth local conformal family of source metrics whose chiral curvatures approach the boundary of definiteness at a point while the source metrics remain regular and the defect energy density has a finite limit. The curvature-reconstructed metric becomes arbitrarily anisotropic at that point.
\end{proposition}
\begin{proof}
Choose round normal coordinates \(y\) at \(p\), and set
\[
f=\tfrac12\bigl((y^1)^2-(y^2)^2\bigr),\qquad
g_\varepsilon=e^{2\varepsilon f}g_0,\qquad z=\varepsilon/K.
\]
At \(p\), \(f=df=0\), so \(g_\varepsilon(p)=g_0(p)\). The conformal Ricci formula gives, exactly at this point,
\[
s=12K,\qquad
\operatorname{Ric}_0=\operatorname{diag}(-2\varepsilon,2\varepsilon,0,0).
\]
The Weyl curvature vanishes identically. In the paired chiral bases
\[
(e^{12}+e^{34},e^{13}-e^{24},e^{14}+e^{23}),\qquad
(e^{12}-e^{34},e^{13}+e^{24},e^{14}-e^{23}),
\]
the curvature blocks are
\[
C=KI,\qquad B=\varepsilon\operatorname{diag}(0,-1,-1).
\]
Consequently,
\[
M=K^2\operatorname{diag}(1,1-z^2,1-z^2),\qquad
Q=\frac{3}{3-2z^2}\operatorname{diag}(1,1-z^2,1-z^2).
\]
The positive range at \(p\) is \(|z|<1\).

To reconstruct the conformal metric, write its determinant-one part as
\(\operatorname{diag}(e^\alpha,e^{-\alpha},1,1)\). Its selfdual graph has entries
\(0,\tanh(\alpha/2),\tanh(\alpha/2)\), and must equal
\(C^{-1}B=\operatorname{diag}(0,-z,-z)\). Therefore
\(e^\alpha=(1-z)/(1+z)\). The volume normalization gives
\begin{equation}\label{eq:degenerate-metric}
\widehat g_A(p)=
\sqrt{1-\frac23z^2}\,
\operatorname{diag}\left(
\frac{1-z}{1+z},\frac{1+z}{1-z},1,1
\right).
\end{equation}
As \(z\to1\), one eigenvalue tends to zero and another to infinity. Meanwhile,
\[
Q\longrightarrow\operatorname{diag}(3,0,0),\qquad
|Q-I|^2=\frac{6z^4}{(3-2z^2)^2},
\]
and hence
\begin{equation}\label{eq:bounded-density}
|Q-I|^2\,\frac{\mu_A}{dV_g}
=\frac{4K^2z^4}{3-2z^2}\longrightarrow4K^2.
\end{equation}
Every expression follows from the exact curvature jet at \(p\), so no limiting numerical argument is involved.
\end{proof}

The proposition is a local family of curvature data realized by metrics, not a solution of the connection flow. A cutoff extends the source family to smooth metrics on the sphere, but it need not preserve definiteness everywhere up to the parameter at which the displayed point degenerates. Thus the example proves a failure of pointwise algebraic coercivity; it does not construct a global singular trajectory or disprove a flow-specific preservation theorem.



\section{The topological consequence of reaching the ASD locus}
\label{sec:ctx-topology}

Let $X$ be a closed smooth manifold homotopy equivalent to $S^4$.
Then $X$ is orientable and simply connected, and
\begin{equation}\label{eq:ctx-sphere-topology}
 \chi(X)=2,\qquad \tau(X)=0,\qquad H^2(X;\mathbb Z)=0.
\end{equation}
The next result separates the geometric existence problem from the
topological conclusion once the required geometry has been obtained.

\begin{theorem}[ASD existence on a smooth homotopy four-sphere]
\label{thm:ctx-ASD-sphere}
For such an $X$, the following statements are equivalent:
\begin{enumerate}
 \item $X$ is diffeomorphic to the standard $S^4$;
 \item $X$ admits a locally conformally flat Riemannian metric;
 \item $X$ admits an ASD Riemannian metric.
\end{enumerate}
Moreover every ASD Riemannian metric on $X$ is conformally
diffeomorphic to the round sphere.
\end{theorem}
\begin{proof}
Pulling back the round metric proves the first implication, and local
conformal flatness implies ASD.  For the converse, the
Chern--Weil signature identity is
\begin{equation}\label{eq:ctx-signature}
 12\pi^2\tau(X)
 =\int_X\bigl(|W^+|^2-|W^-|^2\bigr)\,dV_g.
\end{equation}
If $W^+=0$, equation~\eqref{eq:ctx-sphere-topology} forces
$\int_X|W^-|^2\,dV_g=0$.  Positivity and continuity imply $W^-=0$
pointwise.  In dimension four, $W=0$ is equivalent to local conformal
flatness.

For completeness, the global step is the developing-map argument
underlying Kuiper's theorem~\cite{Kuiper1949}.  Regard each local
conformal chart as taking values in $S^4$, using stereographic
projection.  Liouville's theorem in dimension at least three makes
each transition the restriction of a unique M\"obius transformation.
Continuing one chart along paths therefore defines a conformal local
diffeomorphism from the universal cover to $S^4$; continuation around
homotopic paths agrees.  Since $X$ is simply connected, this gives
$\mathrm{dev}:X\to S^4$.  Its image is open, and compactness makes it
closed, hence all of $S^4$.  A proper local diffeomorphism is a
covering map.  Both spaces are connected and the target is simply
connected, so this covering is a diffeomorphism.  The local conformal
factors agree, giving the asserted conformal diffeomorphism.
\end{proof}

The theorem requires no Einstein equation and no sign condition on
scalar curvature.  It is a classical consequence of the signature
formula and conformal developing theory, not a new classification
theorem.  The unresolved statement is the existence of an ASD metric
on an arbitrary smooth homotopy four-sphere.  By
Theorem~\ref{thm:ctx-ASD-sphere}, proving that statement for all such
$X$ would prove the smooth four-dimensional Poincar\'e conjecture.

\begin{corollary}[Exclusions for vacuum and neutral geometry]
\label{cor:ctx-wave-exclusions}
A smooth homotopy four-sphere admits neither a Ricci-flat ASD
Riemannian metric nor a pseudo-Riemannian metric of signature $(2,2)$.
\end{corollary}
\begin{proof}
In the first case, the preceding signature argument gives $W=0$.
Ricci-flatness then gives the full curvature tensor
$\operatorname{Rm}=0$.  The Euler form of a flat tangent connection
vanishes, contradicting $\chi(X)=2$.

In the second case, a neutral metric supplies a splitting
$TX=E_+\oplus E_-$ into positive and negative rank-two subbundles.
Since $H^1(X;\mathbb Z_2)=0$, both are orientable.  Their Euler
classes lie in $H^2(X;\mathbb Z)=0$, so the Whitney product formula
gives $e(TX)=e(E_+)e(E_-)=0$, again contradicting $\chi(X)=2$.
\end{proof}

Thus neither a Ricci-flat ASD target nor a neutral-signature metric can supply the desired geometry on a smooth homotopy four-sphere. The round Riemannian model retains the nonzero scalar curvature needed for this target.



\section{Outstanding analytical and global challenges}\label{sec:challenges}

The calculations identify a stable linear mechanism and the variables that must be controlled beyond it. They leave the following problems unresolved.

\subsection{A nonlinear equation for the reconstructed geometry}
The exact velocity formula determines \(\widehat g_t\) from \(d_AA_t\), but it is not a closed equation in \(\widehat g\). The compatibility defect
\[
\mathcal T_A=d_A\Sigma_A
=-K^{-1}(d_AQ^{-1/2})\wedge F_A
\]
measures the difference between \(A\) and the chiral Levi--Civita connection of \(\widehat g_A\). A natural next calculation is to derive a coupled evolution for \(\widehat g_A\), \(Q\), and this defect, in a specified diffeomorphism and bundle gauge. Its quadratic expansion should determine exactly which nonlinear terms depart from Ricci flow. The cubic comparison of metric energies does not answer this independent-connection question.

\subsection{Preservation of a quantitative definite region}
A useful continuation estimate would include
\[
\lambda_{\min}(Q)\geq\kappa>0.
\]
Equation~\eqref{eq:min-eigen-defect} explains why a uniform small-defect estimate would suffice locally. The linear heat equation provides such decay for its linearized variable. A nonlinear analogue requires control of the derivative terms and of the moving reconstructed metric. No tensor maximum principle providing this bound is proved here. The finite-density example shows that finiteness of the defect density supplies no pointwise algebraic barrier to degeneration; no nonlinear preservation theorem follows from the energy identity proved here.

\subsection{Regularity, concentration, and geometric compactness}
Even a lower eigenvalue bound for \(Q\) is not a complete regularity statement. One must control the scale and spatial concentration of curvature, the derivatives entering \(\mathcal T_A\), and the geometry of \(\widehat g_A\), modulo diffeomorphisms. Fixed total volume does not imply local noncollapse or a curvature bound. A global argument would need estimates that exclude concentration or describe its possible limits and continuation.

\subsection{The conformal sector and the metric pullback functional}
The metric pullback energy agrees to second order with the Weyl energy, but its conformal kernel at the round metric is not an exact nonlinear conformal invariance. Its unrestricted metric gradient can leave the conformally flat locus. A well-posed nonlinear metric formulation therefore cannot remove every conformal variation as though it were a gauge direction. The connection and metric formulations require separate gauge analysis and separate existence theorems.

\subsection{Initial data and the topology of the four-sphere}
For an arbitrary smooth homotopy four-sphere, this work does not construct a definite connection in the component needed for a positive perfect limit. Nor does it show that every suitable initial connection reaches the round neighborhood. The distinction between an arbitrary definite connection and the chiral Levi--Civita connection of a chosen metric may provide useful flexibility, but it does not remove the existence problem.

If an appropriate flow existed globally, retained definiteness and compactness, and converged smoothly to a perfect connection of the required sign, its reconstructed metric would give the desired geometric conclusion. Each hypothesis contains substantive work. The local Ricci correspondence proved here specifies the behavior near the target and the nonlinear quantities requiring estimates; it does not supply those global hypotheses.

\subsection{Generalized isothermal representations away from integrability}
An extension of the generalized-coordinate idea must retain the obstruction \(W^+\) rather than demand an exact equivalence to a flat integrable twistor distribution. The curvature matrix and compatibility defect provide concrete connection variables near round geometry. What remains is to relate these variables to a canonical almost-twistor or generalized-coordinate construction with suitable domains, transition laws, and reconstruction estimates. Such a construction should preserve the information carried by the conformal scale and trace-free Ricci curvature.

The immediate analytical objective is therefore a nonlinear estimate that couples decay of the curvature defect to quantitative control of the reconstructed metric. The preceding identities isolate that objective without presupposing the global regularity that a four-dimensional uniformization argument would require.



% Bibliography entries follow the order of first citation.
\begin{thebibliography}{99}
\raggedright
\bibitem{Donaldson2011}
S.~K. Donaldson,
\emph{Riemann Surfaces},
Oxford Graduate Texts in Mathematics, vol.~22,
Oxford University Press, Oxford, 2011.
\href{https://doi.org/10.1093/acprof:oso/9780198526391.001.0001}{doi:10.1093/acprof:oso/9780198526391.001.0001}.

\bibitem{Hamilton1982}
R.~S. Hamilton,
\emph{Three-manifolds with positive Ricci curvature},
J. Differential Geom. \textbf{17} (1982), no.~2, 255--306.
\href{https://doi.org/10.4310/jdg/1214436922}{doi:10.4310/jdg/1214436922}.

\bibitem{Perelman2002}
G. Perelman,
\emph{The entropy formula for the Ricci flow and its geometric applications},
arXiv:math/0211159 [math.DG] (2002).
\url{https://arxiv.org/abs/math/0211159}.

\bibitem{Perelman2003Surgery}
G. Perelman,
\emph{Ricci flow with surgery on three-manifolds},
arXiv:math/0303109 [math.DG] (2003).
\url{https://arxiv.org/abs/math/0303109}.

\bibitem{Perelman2003Extinction}
G. Perelman,
\emph{Finite extinction time for the solutions to the Ricci flow on certain three-manifolds},
arXiv:math/0307245 [math.DG] (2003).
\url{https://arxiv.org/abs/math/0307245}.

\bibitem{Smale1961}
S. Smale,
\emph{Generalized Poincar\'e's conjecture in dimensions greater than four},
Ann. of Math. (2) \textbf{74} (1961), no.~2, 391--406.
\href{https://doi.org/10.2307/1970239}{doi:10.2307/1970239}.

\bibitem{Milnor1956}
J. Milnor,
\emph{On manifolds homeomorphic to the 7-sphere},
Ann. of Math. (2) \textbf{64} (1956), no.~2, 399--405.
\href{https://doi.org/10.2307/1969983}{doi:10.2307/1969983}.

\bibitem{Freedman1982}
M.~H. Freedman,
\emph{The topology of four-dimensional manifolds},
J. Differential Geom. \textbf{17} (1982), no.~3, 357--453.
\href{https://doi.org/10.4310/jdg/1214437136}{doi:10.4310/jdg/1214437136}.

\bibitem{FreedmanGompfMorrisonWalker2010}
M.~H. Freedman, R.~E. Gompf, S. Morrison, and K. Walker,
\emph{Man and machine thinking about the smooth 4-dimensional Poincar\'e conjecture},
Quantum Topol. \textbf{1} (2010), no.~2, 171--208.
\href{https://doi.org/10.4171/QT/5}{doi:10.4171/QT/5}.

\bibitem{Fine2011}
J. Fine,
\emph{A gauge theoretic approach to the anti-self-dual Einstein equations},
arXiv:1111.5005v2 [math.DG] (2011).
\url{https://arxiv.org/abs/1111.5005v2}.

\bibitem{AHS1978}
M.~F. Atiyah, N.~J. Hitchin, and I.~M. Singer,
\emph{Self-duality in four-dimensional Riemannian geometry},
Proc. Roy. Soc. London Ser. A \textbf{362} (1978), no.~1711, 425--461.
\href{https://doi.org/10.1098/rspa.1978.0143}{doi:10.1098/rspa.1978.0143}.

\bibitem{Penrose1976}
R. Penrose,
\emph{Nonlinear gravitons and curved twistor theory},
Gen. Relativ. Gravit. \textbf{7} (1976), 31--52.
\href{https://doi.org/10.1007/BF00762011}{doi:10.1007/BF00762011}.

\bibitem{DonaldsonKronheimer1990}
S.~K. Donaldson and P.~B. Kronheimer,
\emph{The Geometry of Four-Manifolds},
Oxford Mathematical Monographs,
Clarendon Press, Oxford, 1990.
\href{https://doi.org/10.1093/oso/9780198535539.001.0001}{doi:10.1093/oso/9780198535539.001.0001}.

\bibitem{Witten1978}
E. Witten,
\emph{An interpretation of classical Yang--Mills theory},
Phys. Lett. B \textbf{77} (1978), 394--398.
\href{https://doi.org/10.1016/0370-2693(78)90585-3}{doi:10.1016/0370-2693(78)90585-3}.

\bibitem{Baker2023}
E.~B. Baker III,
\emph{Generalized conformal maps as classical symmetries of Yang--Mills fields},
arXiv:2301.00029 [math-ph] (2023), and references within.
\url{https://arxiv.org/abs/2301.00029}.

\bibitem{Ward1980}
R.~S. Ward,
\emph{Self-dual space-times with cosmological constant},
Commun. Math. Phys. \textbf{78} (1980), 1--17.
\href{https://doi.org/10.1007/BF01941967}{doi:10.1007/BF01941967}.

\bibitem{FineKrasnovPanov2014}
J. Fine, K. Krasnov, and D. Panov,
\emph{A gauge theoretic approach to Einstein 4-manifolds},
New York J. Math. \textbf{20} (2014), 293--323.
\url{https://nyjm.albany.edu/j/2014/20-18.html}.

\bibitem{GurskyViaclovsky2015}
M.~J. Gursky and J.~A. Viaclovsky,
\emph{Rigidity and stability of Einstein metrics for quadratic curvature functionals},
J. Reine Angew. Math. \textbf{700} (2015), 37--91.
Preprint: arXiv:1105.4648v2 [math.DG].
\url{https://arxiv.org/abs/1105.4648}.

\bibitem{Kuiper1949}
N.~H. Kuiper,
\emph{On conformally-flat spaces in the large},
Ann. of Math. (2) \textbf{50} (1949), no.~4, 916--924.
\href{https://doi.org/10.2307/1969587}{doi:10.2307/1969587}.

\bibitem{Kodaira1962}
K. Kodaira,
\emph{A theorem of completeness of characteristic systems for analytic
families of compact submanifolds of complex manifolds},
Ann. of Math. (2) \textbf{75} (1962), 146--162.
\href{https://doi.org/10.2307/1970424}{doi:10.2307/1970424}.

\bibitem{Kodaira1963}
K. Kodaira,
\emph{On stability of compact submanifolds of complex manifolds},
Amer. J. Math. \textbf{85} (1963), 79--94.
\href{https://doi.org/10.2307/2373187}{doi:10.2307/2373187}.

\bibitem{Shaw1986}
W.~T. Shaw,
\emph{An ambitwistor description of bosonic or supersymmetric minimal
surfaces and strings in four dimensions},
Class. Quantum Grav. \textbf{3} (1986), 753--761.
\href{https://doi.org/10.1088/0264-9381/3/5/006}{doi:10.1088/0264-9381/3/5/006}.

\bibitem{IsenbergYasskinGreen1978}
J. Isenberg, P.~B. Yasskin, and P.~S. Green,
\emph{Non-self-dual gauge fields},
Phys. Lett. B \textbf{78} (1978), 462--464.
\href{https://doi.org/10.1016/0370-2693(78)90486-0}{doi:10.1016/0370-2693(78)90486-0}.

\bibitem{MasonSkinner2006}
L.~J. Mason and D. Skinner,
\emph{An ambitwistor Yang--Mills Lagrangian},
Phys. Lett. B \textbf{636} (2006), 60--67.
\href{https://doi.org/10.1016/j.physletb.2006.02.061}{doi:10.1016/j.physletb.2006.02.061}.
\end{thebibliography}
\end{document}

